%% ============================================================
%% main.tex — Image-Space Commitment Arbitration for Infrared–Visible
%% Image Fusion (Springer sn-jnl, sn-mathphys-num references)
%%
%% Assembled by formatter_agent (Phase 7, 2026-09-14) from
%% writing/00-abstract.md, 01, 02, 04, 05, 06; bibliography from
%% writing/07-references.md (citation-check, dispositions applied).
%% Preamble mirrors the official sn-article.tex v3.1 template.
%% ============================================================

\documentclass[pdflatex,sn-mathphys-num]{sn-jnl}

%%%% Standard Packages
\usepackage{graphicx}%
\usepackage{multirow}%
\usepackage{amsmath,amssymb,amsfonts}%
\usepackage{amsthm}%
\usepackage{mathrsfs}%
\usepackage[title]{appendix}%
\usepackage{xcolor}%
\usepackage{textcomp}%
\usepackage{manyfoot}%
\usepackage{booktabs}%
\usepackage{algorithm}%
\usepackage{algorithmicx}%
\usepackage{algpseudocode}%
\usepackage{listings}%

%% as per the requirement new theorem styles can be included as shown below
\theoremstyle{thmstyleone}%
\newtheorem{theorem}{Theorem}%
\newtheorem{proposition}[theorem]{Proposition}%

\theoremstyle{thmstyletwo}%
\newtheorem{example}{Example}%
\newtheorem{remark}{Remark}%

\theoremstyle{thmstylethree}%
\newtheorem{definition}{Definition}%

\raggedbottom

\begin{document}

\title[Image-Space Commitment Arbitration for IVIF]{Image-Space Commitment Arbitration for Infrared--Visible Image Fusion: A Controlled Falsification and Mechanism-Audit Study}

\author*[1]{\fnm{Xiaotong} \sur{Gao}}%\email{}

\affil*[1]{\orgname{Qingdao University, }\orgaddress{\city{Qingdao}, \country{China}}}

%%--------------------------------------------
%% Abstract (verbatim from writing/00-abstract.md, English abstract)
%%--------------------------------------------
\abstract{At pixels where infrared and visible images disagree, a local convex blend of the two registered source images cancels the structure that committing to the salient modality would preserve --- a pointwise, image-space argument about one operator, not a claim about how any particular network behaves internally. This paper reports a controlled falsification and mechanism audit, not a leaderboard entry. We first test the received diagnosis; it does not replicate: on the selective state-space backbone S4Fusion, modulating the state dynamics in feature space degrades the structural metrics on one seed, improves them on a second, and is reproduced in full by a content-free reliability map, with only QABF's median down. We then build the mechanism that diagnosis implies, Commitment-Gated Arbitration (CGA), which detects cross-modal conflict from image gradients and applies a learned, zero-initialized, image-space per-pixel commitment toward the salient modality, so that the official checkpoint loads exactly. The mechanism is active at inference, but its benefit is not stable: under the budget-matched convention the criterion is met on one seed of three, and a post-hoc cluster-aware re-analysis that treats the image as the unit of observation leaves the same single seed positive, so what survives is a seed-specific, non-replicated effect rather than a replicated gain; significant SF and AG costs accompany it. Content controls show why the trigger cannot be credited: an all-ones map that commits everywhere raises high-conflict recall above the learned map while lowering mAP50 below it, and destroying the map's arrangement changes nothing measurable. What remains is a mechanism implemented and verified but falsified as a source of reliable gain, whose operator is nevertheless not free: no parameter-free rule reproduces the trained selector, which sharpens to a near-binary decision on a sparse support. Training instability, a distillation failure and a pre-freeze cross-dataset arm that reverses direction are reported as negative results; code is released. All results are scoped to a single-testbed, single-detector controlled case.}

\keywords{infrared--visible image fusion, conflict arbitration, selective state-space models, zero-initialized gating, object detection, image quality assessment}

\maketitle

%%=============================================================
%% §1 Introduction
%% Source: writing/01-intro-mechanism.md — Draft paragraphs,
%% "Why S4Fusion as the testbed" items, and "Deliberately not
%% claimed" paragraph; process material (outline, claim–evidence
%% map, assumptions) omitted per assembly rules.
%%=============================================================
\section{Introduction}\label{sec:intro}

Selective state-space models (SSMs) have become a backbone of infrared--visible image fusion because their linear-time global receptive field scales better than self-attention on large inputs \cite{ma2025s4fusion,gu2023mamba}. Their usefulness at high resolution is constrained by state capacity: the recurrent state is bounded, long sequences dilute what each slot retains --- beyond that capacity the model fails even to learn what to forget \cite{chen2025stuffed} --- and the transition is governed by the decay matrix's spectral properties, so poorly managed states drift or vanish \cite{lu2025mambamodulation,ali2025hidden}. Hence the hypothesis that high-resolution fusion failures originate \emph{inside} the state dynamics: that injecting per-pixel boundary reliability into the selective parameters $B$, $C$ and $\Delta$ would re-allocate capacity toward modality boundaries.

We tested it under controlled conditions --- every configuration fine-tuned from the same public S4Fusion weights, same budget, same overlap-tiled inference protocol on M3FD \cite{liu2022tardal} --- and the result is a falsification that does not survive repetition. On seed 42 the state-modulated model significantly \emph{degraded} the structural metrics (SF $-0.019$, AG $-0.013$, QABF $-0.0005$, Holm $p < 0.05$); a second seed meets the same criterion in the \emph{opposite} direction (SF $+0.033$, AG $+0.017$), a third degrades only the QABF median, and the one quantity stable across all three is that median, down by $0.0005$. Under this parameterization feature-space modulation of the state dynamics is therefore neither established nor excluded, and the pivot below rests on the design argument of Section~\ref{sec:methods:problem} rather than on this arm.

Where the two modalities disagree --- opposite edge orientation or polarity at a pixel --- a local convex combination of the two \emph{registered source images} cancels structure, whereas committing to the locally salient modality preserves it; Section~\ref{sec:methods:problem} states the conditions, and notes that this is a statement about a pointwise image-space operator rather than a general property of feature-space fusion. An exploratory training-free check pointed the same way --- committing toward the max-response modality recovered objects that blending loses --- and we record it as a \emph{pilot}, since it predates this study's freeze and has no frozen method record; no result here depends on it. We therefore introduce \textbf{Commitment-Gated Arbitration (CGA)}: it detects per-pixel modal conflict from image gradients, learns a per-pixel commitment toward the salient modality, and applies it through a zero-initialized gate, so that loading the official checkpoint reproduces the baseline exactly \cite{bachlechner2021rezero}. The part of CGA that carried its argument is the restriction of commitment to conflict regions, and Section~\ref{sec:exp:controls} tests it --- with the result that the restriction is \emph{not} supported: an all-ones map raises high-conflict recall above the learned map while lowering mAP50 below it, and permuting the conflict map changes nothing measurable.

S4Fusion is the testbed because it is an open published selective state-space backbone whose released weights make the comparison \emph{controlled}, and because the zero-initialized gate validates the module \emph{against} the weights it modifies. The study is scoped to it; no claim is made that its behaviour typifies scan-based fusion backbones as a class (Section~\ref{sec:limitations}). The mechanism is reported as built and active at inference, and the restriction that motivated it as empirically unsupported; the two claims stay apart throughout.

\emph{Deliberately not claimed}: that S4Fusion is the best fusion model (the comparison table shows EMMA \cite{zhao2024equivariant} leading SF and AG and Diff-IF \cite{yi2024diffif} leading MI, reported as-is); that the choice was mechanistically privileged; that the conflict-gated gain is replicated; or that any result here transfers beyond this backbone, this dataset and this detector. The single-baseline convention clears the high-conflict criterion on all three seeds, but it compares every seed against one shared baseline; the budget-matched convention clears it on one seed of three, and a post-hoc re-analysis that takes the \emph{image} as the unit leaves that one seed alone, so the detection evidence supports a \emph{seed-specific, non-replicated} effect only (Section~\ref{sec:exp:verdicts}). Neither convention reaches the stricter thresholds.

Four things, none of them a performance claim. \emph{First, a zero-initialized image-space adapter compatible with the weights it modifies}: a per-pixel source-selection network combined with the backbone's fused image through one zero-initialized scalar gate, so the augmented model reproduces the official checkpoint exactly at load time and its effect is measured only after training opens the gate (Sections~\ref{sec:methods:cga}--\ref{sec:methods:training}). \emph{Second, a falsification that does not itself replicate}: the intervention the literature's capacity diagnosis implies fails the structural criterion on the seed it was run on, meets the same criterion in the opposite direction on a second, and is reproduced in full by a content-free reliability map, so its effect is not a property of boundary reliability at all (Section~\ref{sec:exp:verdicts}). \emph{Third, the separation of the trained selector from the parameter-free rules it resembles}: the benefit does not replicate, but the operator is not free --- seven parameter-free rules on two supports are each worse than the learned selector at the detection endpoint, and the sharpest, the same architecture with only the selector reset, is slightly worse than doing nothing (Table~\ref{tab:free}). \emph{Fourth, a protocol}: criterion layering, both base conventions reported as non-comparable, replication counted per seed, and a post-hoc cluster-aware re-analysis that treats the image rather than the object as the unit of observation (Sections~\ref{sec:methods:design} and~\ref{sec:exp:verdicts}). The protocol is what turns a weak result into a legible one, and it is the part of this work most directly reusable.

%%=============================================================
%% §2 Related work
%% Source: writing/02-related-work.md — six family sections;
%% classification map converted to transitional prose
%% (journal style); Notes section omitted.
%%=============================================================
\section{Related work}\label{sec:related}

CNN fusion builds on autoencoder and dense designs (DenseFuse \cite{li2019densefuse}, FusionDN \cite{xu2020fusiondn}, U2Fusion \cite{xu2022u2fusion}) and on adversarial generation (FusionGAN \cite{ma2019fusiongan}); Transformer fusion added global context (CDDFuse \cite{zhao2023cddfuse}, MetaFusion \cite{zhao2023metafusion}, EMMA \cite{zhao2024equivariant}), and diffusion fusion casts the problem as iterative denoising (DDFM \cite{zhao2023ddfm}, Diff-IF \cite{yi2024diffif}). Across all three families the operator remains a weighted combination of modality evidence: at a pixel where the sources disagree the network chooses how much of each to mix rather than deciding which to trust.

Selective state spaces entered IVIF nearly simultaneously --- S4Fusion \cite{ma2025s4fusion}, the wavelet-domain W-Mamba \cite{zhang2025wmamba} and the multi-modal FusionMamba \cite{xie2024fusionmamba} --- inheriting linear-time global context from the SSM family \cite{gu2023mamba,dao2024transformers}, and their limitation is state capacity: an overloaded recurrent state loses its ability to manage what it retains \cite{chen2025stuffed}, and only the A/$\Delta$ parameters govern history decay \cite{ali2025hidden,lu2025mambamodulation}. In the implementations we reviewed they scan and blend rather than select, and the 2026 wave modifies \emph{scan topology} \cite{zhu2026interactive,cao2026shuffle,wang2026selfsupervised}, i.e.\ \emph{where the scan looks}. Our own ablation (Section~\ref{sec:methods:design}) finds that modulating the state dynamics with boundary reliability does not improve structural metrics under its parameterization, which is why what we test instead is a selection decision at conflicting pixels.

Task-guided fusion --- SeAFusion \cite{tang2022seafusion}, SegMiF/M2Fusion \cite{liu2023segmif}, MetaFusion \cite{zhao2023metafusion} and lightweight KD variants \cite{zhao2023lightweight,wang2025semantic} --- guides \emph{where} to look while still weighting features at conflicting pixels and needs task labels at inference. Boundary and uncertainty cues are an active direction (BUFNet \cite{zhou2026bufnet}, ShadowMamba \cite{zhu2026shadowmamba}, edge-prior-guided IVIF \cite{yuan2025segment}, edge-guided dual attention \cite{zou2023lightweight}, uncertainty-aware detection fusion \cite{zhao2024uncertainty}).

Closest to our setting, ARC3Fusion \cite{tian2026adaptive} models consensus--complementarity--\emph{conflict} in IVIF and GatedFusion-Net \cite{brenner2026gatedfusion} computes per-pixel modality weights, but both are \emph{continuous, feature-space} coordinators: three-branch Softmax weights summing to one, and per-pixel Sigmoid masks. Image-space per-pixel \emph{hard} selection has a separate lineage --- the classical select/choose rules \cite{burt1993enhanced,zhang1999categorization,piella2003general} and the learning-based multi-focus decision-map line \cite{ma2021endtoend} --- but same-modality and defocus-driven. Learned \emph{discrete} routing has also entered IVIF: a concurrent preprint \cite{ma2026dynamic} selects among three experts per pixel via a Gumbel-Softmax straight-through estimator, the closest existing instance of per-pixel hard choice in cross-modal fusion, though it operates on feature components and routes on learned informativeness rather than cross-modal gradient conflict. To the best of our knowledge no prior IVIF method converts cross-modal gradient-orientation conflict into a learned, image-space, per-pixel hard selection. That is a statement of absence in the literature we reviewed, not a proof of novelty, and the characterizations differ in verification depth: the two published neighbours were read in full, while the concurrent preprint was verified at abstract level only, so nothing is asserted about its unstated components. Table~\ref{tab:positioning} states the resulting distinction on the axes that can be discriminated --- operating space, selector discreteness, final-output discreteness, support, supervision, trigger quantity and evidence status --- and it is that axis distinction, not the difference magnitude, that this paper claims. The last row of the table is the one the rest of the paper tests, and under the matched-budget convention it is only partly supported: the operating space and the zero-initialized compatibility hold, the final output is a continuous residual arbitration rather than a discrete pick, and the trigger does not (Section~\ref{sec:exp:controls}).

\begin{table*}[t]
\caption{Mechanism axes for conflict-aware fusion, stated at the granularity at which the selection operator can actually be discriminated. ``Conc.'' marks concurrent work. The evidence column of the last row records how far each characterization was verified: full-text reading for the two published neighbours, abstract-level verification only for the concurrent preprint, and the classical lineage from its own published statements. The final row is the mechanism this paper builds and tests.}
\label{tab:positioning}
\scriptsize
\setlength{\tabcolsep}{3pt}
\resizebox{\textwidth}{!}{%
\begin{tabular}{@{}p{2.15cm}p{2.5cm}p{2.55cm}p{2.4cm}p{2.6cm}p{2.3cm}@{}}
\toprule
Axis & ARC3Fusion & GatedFusion-Net & Ma et al.\ (conc.) & select/choose lineage & \textbf{CGA (ours)} \\
\midrule
Operating space & feature & feature & feature & image & \textbf{image} \\
Selector discreteness & continuous: three-branch Softmax summing to one & continuous: per-pixel Sigmoid masks & discrete: Gumbel--Softmax straight-through & discrete: binary choose & \textbf{continuous by construction, near-binary after training} \\
Final-output discreteness & continuous & continuous & not asserted (abstract-level) & discrete & \textbf{continuous residual arbitration} \\
Support & all positions & all pixels & per-pixel expert set & disagreement / defocus regions & \textbf{learned sparse map ($\overline{M}=0.0447$)} \\
Supervision & end-to-end task loss & end-to-end task loss & end-to-end task loss & rule, none & \textbf{self-supervised conflict prior + commit loss; no detection labels} \\
Trigger quantity & edge strength \& orientation mismatch (explicit) & learned modality confidence (implicit) & learned expert informativeness (implicit) & local sharpness & \textbf{cross-modal gradient orientation \& polarity} \\
Detector endpoint reported & no & no (segmentation) & no & no & \textbf{yes: recall and mAP50} \\
Evidence status & full text & full text & \textbf{abstract only} & published statements & \textbf{this paper, full analysis} \\
\botrule
\end{tabular}}%
\end{table*}

%%=============================================================
%% §3 Methods
%% Source: writing/04-methods.md — five subsections;
%% (gate types, boundary modes, CGA); assembly markers retained.
%% placeholder comment replaced by the measured efficiency table (B-2, 2026-09-16)
%% comments with qualitative in-text wording.
%%=============================================================
\section{Methods}\label{sec:methods}

\subsection{Problem formulation and design requirements}\label{sec:methods:problem}

Let $I_{\mathrm{ir}}, I_{\mathrm{vi}} \in [0,1]^{1\times H\times W}$ be a fixed infrared--visible pair, the visible image reduced to its luminance (Y) channel. The backbone, S4Fusion \cite{ma2025s4fusion}, encodes both modalities with selective state-space (SSM) encoders, fuses them at multiple scales in bidirectional-scan fusion blocks, and decodes a fused luminance image $Y = f_\theta(I_{\mathrm{ir}}, I_{\mathrm{vi}})$. The question this section formalizes is where, in such a pipeline, the decision should be made when the two modalities disagree about the same physical structure.

We define cross-modal conflict at a pixel as the conjunction of two conditions: both modalities carry structure there, and their gradient vectors disagree in orientation or polarity --- the common case at modality boundaries, where a thermal edge and its visible counterpart often have opposite polarity. The argument is made under three explicit conditions: the sources are \emph{registered}, the comparison is \emph{local}, and the operator is a \emph{pointwise convex combination of the two source images} rather than of any intermediate feature tensor. Under those conditions a combination is structurally destructive at a conflicting pixel: if locally $\nabla I_{\mathrm{vi}} \approx -\nabla I_{\mathrm{ir}}$ with common magnitude $g$, then $F = w\,I_{\mathrm{ir}} + (1-w)\,I_{\mathrm{vi}}$ has gradient magnitude $|\nabla F| \approx |2w-1|\,g$, which vanishes at $w = 1/2$ and is preserved only at the endpoints. For arbitrary magnitudes and orientation agreement $c \in [-1,1]$ the squared magnitude is $\|\nabla F\|^2 = w^2 g_{\mathrm{ir}}^2 + (1-w)^2 g_{\mathrm{vi}}^2 + 2w(1-w) g_{\mathrm{ir}} g_{\mathrm{vi}} c$, whose minimum over $w \in [0,1]$ is $\dfrac{g_{\mathrm{ir}}^2 g_{\mathrm{vi}}^2 (1-c^2)}{g_{\mathrm{ir}}^2 + g_{\mathrm{vi}}^2 - 2 g_{\mathrm{ir}} g_{\mathrm{vi}} c}$ when the interior minimiser is attainable; the anti-parallel equal-magnitude case is its zero limit, and even orthogonal equal-magnitude gradients ($c = 0$) lose half their energy under an equal blend. The prior $\kappa$ of Section~\ref{sec:methods:cga} is defined over the full range of $c$ for that reason. A pointwise image-space weighting therefore cannot preserve a conflicting boundary; it can only choose how much to cancel (Fig.~\ref{fig:soft-vs-hard}). Whether the same conclusion extends to \emph{feature-space} fusion is a separate question and is treated here as a hypothesis rather than a result, since a fusion network applies nonlinear mappings and decoding after any feature-space weighting, so its output is not a convex combination of the inputs; the feature-space arm of Section~\ref{sec:methods:design} is the empirical test of that extension.

\begin{figure}[t]
\centering
\includegraphics[width=0.72\textwidth]{figs/soft_vs_hard.pdf}
\caption{Structure preserved by a blend at a conflicting pixel, as a function of the weight $w$ given to the infrared modality, normalized by the common gradient magnitude $g$. The closed form $|2w-1|$ (line) is reproduced pointwise by a numerical measurement on a synthetic opposite-polarity pair (circles; maximum deviation $6\times10^{-15}$, i.e.\ at floating-point precision). A soft weight reaches zero at $w=1/2$ and cancels the structure entirely; a hard commitment sits at an endpoint and preserves all of it. This is the argument, not the demonstration, that motivates an image-space commitment; whether the commitment has to be \emph{confined} to conflict pixels is a separate question that Section~\ref{sec:exp:controls} answers in the negative.}
\label{fig:soft-vs-hard}
\end{figure}

The natural alternative --- giving the operator more capacity to represent boundaries --- we tested directly: a same-budget ablation degraded the structural metrics on the seed it was run on, but a two-seed repetition meets the same criterion in the opposite direction on one of them (verdict in Section~\ref{sec:exp}). Under this parameterization the capacity-dilution account is consequently neither transferred to nor excluded from structural gains, and the alternative we test next is argued on design grounds rather than on this arm.

These observations induce three design requirements. (R1)~\emph{Image-space operation}: which modality to keep is observable in image space, the state-dynamics arm supplies no replicated evidence that a feature-space intervention transfers, and image-space operation leaves the backbone untouched. (R2)~\emph{Conflict-triggered selectivity}: the mechanism should act only where the modalities conflict, so non-conflict pixels retain the pretrained behaviour and any measured difference is attributable to selection rather than to a global perturbation. (R3)~\emph{Per-pixel commitment behind a zero-initialized gate}: at a conflict pixel the output should move toward the locally salient modality rather than a blend of both, and the mechanism must start as an exact identity so that pretrained weights remain valid. (R2) is the requirement this paper ends up testing rather than relying on: the controls of Section~\ref{sec:exp:controls} do not support it, and we report that failure separately instead of restating the requirement as an accomplishment. (R1) and (R3) hold as stated.

\subsection{Commitment-Gated Arbitration}\label{sec:methods:cga}

Commitment-Gated Arbitration (CGA) implements the three requirements as one image-space module: detect conflict, form a per-pixel commitment toward the salient modality, and apply that commitment through a zero-initialized gate. The gate is designed to act selectively through a learned support map, and whether this conflict-specific localization is necessary is not assumed but tested: the uniform and shuffled-map controls of Section~\ref{sec:exp:controls} do not support it, so the localization is reported here as a design motivation under test rather than as an established property. It is attached once, at the output of the S4Fusion decoder: the augmented forward pass retains the original inputs at native resolution and applies CGA to $(I_{\mathrm{ir}}, I_{\mathrm{vi}}, Y)$, returning the arbitrated image $F$; encoders, fusion blocks and decoder are untouched, and the interface (two source images plus a fused image) is backbone-agnostic in principle. Figure~\ref{fig:cga-module} gives the resulting data path: the priors of Eqs.~(\ref{eq:gm})--(\ref{eq:kappa}) feed the conflict head, the selector forms the commitment, and both are multiplied into a single per-pixel weight before touching the fused image.

\begin{figure}[t]
\centering
\includegraphics[width=0.98\textwidth]{figs/cga_module.pdf}
\caption{The CGA module. Fixed Sobel gradients of the source pair give the orientation-agreement score $c$ and the conflict prior $\hat\kappa$ (Eqs.~\ref{eq:gm}--\ref{eq:kappa}); a learned head turns them into the conflict map $M$, and a learned selector produces the per-pixel commitment $K$. Both act on the backbone's fused output $Y$ only through the single weight $w = \tanh(\alpha)M$, and both $\alpha$ and the selector's final layer are zero-initialized, so the whole module is an exact identity at load time and the frozen S4Fusion parameters are never re-parameterised. Whether $w$ has to be confined to conflict pixels is the question Section~\ref{sec:exp:controls} answers in the negative.}
\label{fig:cga-module}
\end{figure}

\paragraph{Conflict signal.}
For each modality $m \in \{\mathrm{ir}, \mathrm{vi}\}$, fixed $3\times 3$ Sobel kernels scaled by $1/4$ give gradient components $G_x^m, G_y^m$ and magnitude
\begin{equation}\label{eq:gm}
g_m = \sqrt{(G_x^m)^2 + (G_y^m)^2 + \varepsilon}, \qquad \varepsilon = 10^{-6},
\end{equation}
and orientation agreement between the two gradient fields is their cosine
\begin{equation}\label{eq:cos}
c = \frac{G_x^{\mathrm{ir}} G_x^{\mathrm{vi}} + G_y^{\mathrm{ir}} G_y^{\mathrm{vi}}}{g_{\mathrm{ir}}\, g_{\mathrm{vi}} + \varepsilon} \in [-1, 1].
\end{equation}
The conflict prior is the closed form
\begin{equation}\label{eq:kappa}
\kappa = \min\!\left(g_{\mathrm{ir}}, g_{\mathrm{vi}}\right)\, \frac{1 - c}{2}.
\end{equation}
The minimum makes the signal symmetric: conflict requires structure in \emph{both} modalities, so a pixel where only one modality carries an edge contributes a value at the numerical floor rather than exactly zero --- with $\varepsilon = 10^{-6}$ the magnitude of an edgeless modality is $g_m = \sqrt{\varepsilon} = 10^{-3}$, so such a pixel yields $\kappa \approx \tfrac{1}{2}\sqrt{\varepsilon} = 5\times10^{-4}$; measured over 40 M3FD test pairs this floor averages $1.47\times10^{-3}$ of the per-image maximum of $\kappa$, i.e.\ it is negligible in practice but not identically zero, and the phrase ``contributes zero'' is used below in that approximate sense only. Meanwhile $(1-c)/2$ maps aligned gradients ($c \approx 1$) to $\kappa \approx 0$ and opposite polarity ($c \approx -1$) to its maximum. The prior is computed in fp32 even under mixed precision, so that the square-root epsilon remains effective, and is normalized per image by its spatial maximum and clipped, $\hat\kappa = \mathrm{clip}\big(\kappa / (\max_u \kappa + \epsilon);\, 0, 1\big)$, serving as a self-supervised target below. No annotation enters any of these expressions; the signal is derived from the input pair alone. The orientation-agreement factor $(1-c)/2$ is the same quantity that edge-preservation indices weight when they ask how much of an edge's orientation survives fusion \cite{xydeas2000objective}; here it enters a per-pixel gate rather than an aggregate index. We use five symbols with distinct referents throughout, and we do not describe the architecture as a discrete router: $s$ is the \emph{source-selection weight} (continuous, per pixel, from a sigmoid); $K = s\,I_{\mathrm{ir}}+(1-s)\,I_{\mathrm{vi}}$ is the \emph{source mixture}, an attractor toward one source whose value depends on how far $s$ has been driven; $M$ is the learned \emph{support} (intensity) map that decides where the mechanism is allowed to act; $w = \tanh(\alpha)M$ is the signed per-pixel \emph{residual strength}; and $F = Y + w\odot(K-Y)$ is the \emph{final continuously arbitrated output}. The operator is therefore a continuous residual arbitration driven by a selector, not a discrete hard router. Whether the trained selector behaves like a binary choice is an \emph{empirical} property of the trained weights --- we call it a \emph{trained near-binary selector} only where Section~\ref{sec:exp:vis} has measured it --- and ``commitment'' names the operation of moving the output toward $K$ rather than averaging the two sources.

\paragraph{Learned conflict map.}
The prior is pointwise and fixed; the mechanism learns its conflict map from richer evidence. A convolutional head $h_c$ (a $3\times 3$ convolution from 5 to 16 channels, SiLU, a depthwise $3\times 3$ convolution, SiLU, a $1\times 1$ convolution to one channel) receives the stack $[I_{\mathrm{ir}}, I_{\mathrm{vi}}, g_{\mathrm{ir}}, g_{\mathrm{vi}}, \hat\kappa]$ and produces
\begin{equation}\label{eq:M}
M = \sigma\big(h_c([I_{\mathrm{ir}}, I_{\mathrm{vi}}, g_{\mathrm{ir}}, g_{\mathrm{vi}}, \hat\kappa])\big) \in [0,1]^{1\times H\times W},
\end{equation}
a map that can absorb context the pointwise prior cannot --- textured versus structural disagreement, for instance --- while remaining grounded on $\hat\kappa$ through the loss below.

\paragraph{Commit selection.}
The commitment is a per-pixel convex selection between the two source images,
\begin{equation}\label{eq:commit}
K = s\, I_{\mathrm{ir}} + (1-s)\, I_{\mathrm{vi}}, \qquad s = \sigma\big(h_s([I_{\mathrm{ir}}, I_{\mathrm{vi}}, g_{\mathrm{ir}}, g_{\mathrm{vi}}]) + 2p\big),
\end{equation}
where $h_s$ is a $3\times 3$ convolution (4 to 16 channels), SiLU, $1\times 1$ convolution stack and $p = 2\cdot\mathbb{1}[g_{\mathrm{ir}} > g_{\mathrm{vi}}] - 1 \in \{-1, +1\}$ names the modality with the larger local gradient. The final layer of $h_s$ is initialized to zero, so before training the selection reduces to $s = \sigma(2p) \approx 0.88$ toward that modality --- the max-response commit identified by the training-free pre-check reported in the Introduction --- and training may sharpen or override it.

\paragraph{Gate and combination.} A single scalar $\alpha$, zero-initialized, bounds the mechanism globally; the learned conflict map localizes it:
\begin{equation}\label{eq:gate}
w = \tanh(\alpha)\cdot M, \qquad F = Y + w \odot (K - Y),
\end{equation}
equivalently $F = (1-w) \odot Y + w \odot K$ with a per-pixel, \emph{signed} $w \in [-\tanh|\alpha|, +\tanh|\alpha|]$. The sign matters: since $M \ge 0$, a negative $\alpha$ makes $w \le 0$ and $F$ moves \emph{away from} $K$, so the convex reading with $w \ge 0$ holds only for $\alpha \ge 0$ and is used below only in that case. In every checkpoint reported here $\alpha$ is positive and the gate is open --- $\alpha = +4.665$, $+3.674$ and $+3.749$ for seeds 42, 123 and 3407, i.e.\ $\tanh(\alpha) = 0.99982$, $0.99871$ and $0.99889$, with last-epoch values $+4.767$, $+4.551$, $+4.638$ --- so the sign is positive throughout, while the general signed form is retained because the objective does not constrain it. Where $M \to 0$ the output equals $Y$ exactly; where the mechanism acts it moves toward the source mixture by at most $|\tanh(\alpha)|$. At initialization $\tanh(0) = 0$, hence $F \equiv Y$ at every pixel.

\paragraph{Training objective.}
The backbone keeps the repository's standard fine-tuning objective --- an intensity term ($\mathcal{L}_1$ to the pixelwise maximum of the inputs), a gradient term ($\mathcal{L}_1$ to the pixelwise maximum of the Sobel magnitudes), an edge-aware smoothness term, and a cross-resolution scale-consistency term --- at repository defaults. CGA adds two terms. The conflict grounding term is a binary cross-entropy with logits between the head's logit and the detached target $\hat\kappa$ (numerically safe under mixed precision). The conflict-gated commit term is
\begin{equation}\label{eq:commitloss}
\mathcal{L}_{\mathrm{commit}} = \frac{1}{HW}\sum_{u} \bar{M}(u)\,\big(F(u) - K(u)\big)^2, \qquad \bar{M} = \mathrm{detach}(M),
\end{equation}
which pushes the output toward the source mixture in predicted conflict regions only. Since $F - K = (1 - w)(Y - K)$, this term can be reduced along at least three distinct paths: opening the gate (raising $\alpha$ or $M$), moving $K$ toward $Y$ by flattening the selector toward the backbone's own blend, or sharpening the selector so that $K$ and $Y$ happen to coincide where $M$ is large. It is therefore \emph{not} identifiable as evidence that the selector has learned to pick the locally salient modality, and we use it only in its design capacity --- as the term that drives the zero-initialized gate away from zero --- while the claims about what the selector learned rest on the measured selector distribution of Section~\ref{sec:exp:vis} and on the parameter-free counterfactuals of Table~\ref{tab:free}, not on this loss. The total objective is
\begin{equation}\label{eq:total}
\mathcal{L} = \mathcal{L}_{\mathrm{fusion}} + \lambda_c\, \mathcal{L}_{\mathrm{conflict}} + \lambda_k\, \mathcal{L}_{\mathrm{commit}},
\end{equation}
with $\lambda_c = 0.1$ and $\lambda_k = 10.0$ in all reported runs.

\subsection{Zero-initialized compatibility and training protocol}\label{sec:methods:training}

\paragraph{Exact compatibility at load time.}
The gate is the only path through which CGA can modify the output, and it is closed at initialization: $\tanh(0) = 0$ makes $w \equiv 0$ and $F \equiv Y$, regardless of the randomly initialized heads behind the closed gate. Loading the official S4Fusion checkpoint into the augmented architecture therefore reproduces the baseline exactly: every pretrained key loads unchanged, and the only keys absent from the checkpoint are those of the CGA module, all of which act behind the closed gate. This is the ReZero-style identity-at-initialization principle \cite{bachlechner2021rezero} --- a strictly zero-initialized residual path is an exact no-op at load time, unlike small-nonzero schemes such as LayerScale \cite{touvron2021going} --- and it is what makes the study controlled: all arms of Section~\ref{sec:methods:design} start from the same public weights, so any measured difference is attributable to the intervention rather than to architecture surgery. After training, the gate saturates near unity in every reported checkpoint --- $\tanh(\alpha) \in [0.9987, 0.9999]$ across the three seeds (Section~\ref{sec:methods:cga}) --- which excludes the self-falsification scenario in which the gate stays at zero and the measured differences have some other origin.

\paragraph{Training protocol.}
Every arm fine-tunes from the official public checkpoint under a unified schedule: $128\times 128$ random crops, batch size 1 with 8-step gradient accumulation, 2,048 crops per epoch drawn by dataset-balanced sampling from a multi-source training manifest, AdamW with learning rate $2\times 10^{-5}$ and weight decay $10^{-4}$, gradient-norm clipping at 10.0, and mixed-precision training on a 4\,GB-class consumer GPU. The CGA arm extends the schedule from 20 to 32 epochs and assigns the CGA parameters a separate learning rate of $0.01$ with zero weight decay, so that the zero-initialized mechanism can open within a short fine-tuning budget while the backbone is preserved at the base rate. Three seeds (42, 123, 3407) are run for every reported cell; checkpoints are selected by lowest validation loss and written atomically. No detection supervision enters training anywhere: the commit signal is self-derived from the input pair, and the detection endpoint of Section~\ref{sec:methods:design} is evaluation-only --- the detector is never in the training loop.

Two disclosures about this recipe are load-bearing for how every number in Section~\ref{sec:exp} should be read. First, the arm is a \emph{bundle}: module, conflict-grounding loss, commit-weighted loss, a separate high learning rate for the module, a schedule extended from 20 to 32 epochs and a best-validation checkpoint rule are introduced together, and no reported comparison separates them. Closing that gap requires an orthogonal-factor ablation under one budget and one checkpoint rule; that ablation was \textbf{not run}, because the gate authorizing it --- replication of the primary endpoint under a cluster-aware analysis --- was not met (Section~\ref{sec:exp:verdicts}). The consequence is accepted rather than argued away: the reported CGA$-$B0 deltas are the effect of the \emph{recipe as a whole}, and this paper does not claim to have identified the \emph{operator} as their cause. The pre-freeze implementation is labelled \texttt{CGA-v1} wherever that distinction matters. Second, the training run is \emph{not} numerically stable: late-epoch validation NaN appears in all three CGA-line seeds while the baseline line, including the 32-epoch matched control, completes without NaN, so every reported CGA number comes from a best-validation checkpoint selected \emph{before} NaN onset (Section~\ref{sec:exp:negative}). The method is recorded as an unstable experimental prototype rather than a finished pipeline.

\paragraph{Provenance of the pre-registered criterion.} The detection criterion was registered on 9 September 2026, whereas the seed-42 checkpoint behind the main table's CGA row was written on 25 August 2026 and the gate thresholds, the two base conventions and the wording rules were fixed on 13--15 September. The pre-registration is therefore a \textbf{later-stage prospective re-test with a staged lock}: it fixes the criterion and the reporting rules before the budget-matched re-test they govern, and does not retroactively cover the architecture, recipe or checkpoint that preceded it. Section~\ref{sec:exp:repro} gives the item-by-item timeline.

\paragraph{Baseline arm.}
The same-budget baseline (B0) retrains the unmodified architecture from the same checkpoint under the unified 20-epoch schedule with all optional branches disabled, anchoring comparisons to the training regime rather than to published numbers. The 32-epoch extension of the CGA arm was fixed together with its success criteria before the re-test was run, with every other knob held identical; the budget difference to the 20-epoch B0 arm is therefore disclosed rather than claimed away, and the reported CGA$-$B0 deltas are read as the effect of the CGA recipe (gate, commit-weighted loss, gate learning rate, extended schedule) as a whole. Best-validation checkpoints were selected before any NaN onset in every CGA seed (epochs 11/12/26 for s123/s3407/s42, respectively), and the fact that the seed-42 checkpoint was selected at epoch 26 --- beyond the 20-epoch baseline schedule, and adjacent to the NaN that begins at epoch 27 --- is part of the instability disclosure above rather than an incidental detail. A late-training validation-loss instability observed in all three CGA-line seeds, with best checkpoints selected before its onset by the best-validation rule, is reported with its record in Section~\ref{sec:exp}.

\subsection{Falsification design}\label{sec:methods:design}

\paragraph{Design principle.}
The methodological claim of this paper rests on an ordering: hypotheses, pass criteria, and even the narrative disposition for each outcome tier were fixed before the experiments, and the verdicts are recorded. Two hypothesis tracks share the public starting checkpoint. The first (H1) tests the mechanism hypothesis --- that boundary failures of SSM fusion originate in the state dynamics and can be repaired by feature-space modulation. The second (H5) tests the arbitration hypothesis --- that a conflict-gated, image-space commitment produces a detectable downstream gain. The tracks are read independently: a negative verdict on H1 demotes the state-capacity account without touching H5, and vice versa; neither can compensate for the other.

\paragraph{H1 arm: feature-space state modulation.}
Each fusion block is augmented with a boundary-reliability estimator, a small convolutional network over eight feature-structure channels (normalized Sobel magnitudes and absolute Laplacian responses per modality, their edge agreement, per-modality local contrast, and the edge-magnitude difference) plus a pooled two-channel global context, producing a three-way (IR / VI / shared) reliability distribution at every position; the estimator is initialized neutrally (uniform reliability) with weak fixed structural priors that the learned residual may override. Three tanh-bounded, zero-initialized scalars per block then modulate the selective scan. Writing the reliability maps $r_{\mathrm{ir}}, r_{\mathrm{vi}}, r_{\mathrm{sh}}$ and $s_i = \tanh(\beta_i)$ with $\beta$ zero-initialized,
\begin{align}
\gamma_{\mathrm{ir}} &= 1 + s_0\,(r_{\mathrm{ir}} - r_{\mathrm{vi}}), \qquad \gamma_{\mathrm{sh}} = 1 + s_1\, r_{\mathrm{sh}}, \qquad \gamma_{\Delta,m} = 1 + s_2\,(1 - r_m), \nonumber\\
B_m &\leftarrow \gamma_m\, B_m, \qquad C_m \leftarrow \gamma_m\, \gamma_{\mathrm{sh}}\, C_m, \qquad \Delta_m \leftarrow \gamma_{\Delta,m}\, \Delta_m,
\label{eq:h1mod}
\end{align}
with the shared gain applied to a cross-modal component as well. The pass criterion requires improvement in at least two structural metrics (SF, AG, QABF or EOR) under per-image paired Wilcoxon tests with Holm correction at $p < 0.05$; the verdict is a falsification, reported in Section~\ref{sec:exp}.

\paragraph{H5 endpoint.}
The primary endpoint of the arbitration hypothesis is deliberately mechanism-external: detector recall on high-conflict objects of the full 300-image M3FD test split, with an off-the-shelf detector held fixed across all arms (configuration in Section~\ref{sec:exp}). An object's conflict score is the mean of the closed-form prior $\kappa$ (Section~\ref{sec:methods:cga}) over its ground-truth box, computed from the input pair alone; objects strictly above the 66th percentile of conflict scores form the high-conflict subset. The comparison is paired: recovered objects (missed by the baseline, detected with CGA) and lost objects (the converse) form the discordant pairs, and the pass criterion is $p < 0.05$ together with recovered $-$ lost $\geq 3$, with $p$ read from the test appropriate to the analysis layer in which it is evaluated. Because the subset is defined by the same statistic that triggers the mechanism, two guards against selection bias are built in: the endpoint is a downstream quantity (detector recall, not any internal of the fusion), and the stricter gates below add all-object mAP50 and structural non-inferiority.

\paragraph{Statistical units and analysis families.}
The design has a three-level structure that the primary inference has to respect: objects are nested in images, and images are reused across the three training seeds. An object-level exact binomial test treats the discordant objects as independent observations, which they are not --- objects inside one image share that image's fused output and its detection error --- so the analysis is layered explicitly. The original protocol specified an object-level paired endpoint, and that endpoint and its criterion were fixed before the arbitration re-test. After the dependence structure was identified, a post-hoc image-clustered re-analysis was added as a robustness analysis and is reported separately from the pre-specified endpoint. Because multiple objects share the same image-level fused output and detection error, this revision reads the post-hoc image-clustered analysis as the primary interpretation of dependence-aware evidence, and retains the original object-level exact binomial as a sensitivity analysis that is not used alone to support replication. The three layers are named once here and used with those names throughout:

\begin{itemize}
\item \textbf{Post-hoc cluster-aware robustness analysis} --- the reading the detection verdict rests on, on the fixed 300-image test set (Table~\ref{tab:cluster}). The image is the unit of observation. For each seed we report (i) an image-level net-sign test on the per-image difference in recovered-minus-lost high-conflict objects, (ii) an exact cluster sign-flip test on the image scores, and (iii) a percentile bootstrap for the recall difference. The bootstrap resamples the $181$ images that carry the subset with replacement ($B = 2000$, fixed seed), but holds the denominator fixed at the $878$ observed high-conflict objects; its estimand is therefore the \emph{object-weighted} recall difference of the two fixed arms over the fixed image set, not an image-average recall difference and not a population quantity. Fixing the denominator is deliberate: the arms are compared on one and the same already-annotated object set, so letting the resampled object count enter the denominator would inject a between-image object-count variability that this endpoint does not have. A seed counts as supporting the hypothesis only if the image-level test rejects at $0.05$ \emph{and} the bootstrap interval excludes zero; the criterion of $p<0.05$ with recovered $-$ lost $\geq 3$ is then read \emph{within} that layer. This analysis was added after the dependence structure was identified --- that is, after the re-test outcome it describes was already in hand --- so it is a post-hoc robustness analysis of stored predictions and not a pre-registered test; Table~\ref{tab:timeline} records its date on that footing.
\item \textbf{Pre-specified sensitivity endpoint (object level).} The exact binomial (McNemar-style) test on the paired recovered/lost objects is retained and reported beside every cluster-aware result, because the protocol fixed it before the cluster structure was modelled. It is a sensitivity analysis: it can show that a difference exists in aggregate, but it cannot support a replication claim on its own, and the paper does not use it as one. The two object-level tables are read only in this capacity --- Table~\ref{tab:h5} under the single-baseline convention and Table~\ref{tab:h5e32} under the budget-matched one --- as sensitivity material and as a record of the original protocol's own endpoint. Neither carries the verdict; that is Table~\ref{tab:cluster}.
\item \textbf{Descriptive pooled analysis.} The pooled across-seed statistic and the pooled $n=900$ quality analysis describe three already-trained models on one fixed image set. They are neither independent training replications nor a basis for extrapolating to the population of training runs, and they are reported in that capacity only.
\end{itemize}

No equivalence claim is made anywhere: where two arms are not distinguished by a test, the wording is that \emph{no difference was detected} at the reported power, not that the arms are equivalent. This applies in particular to the VIF comparison, to the shuffled-map control, and to the gate-C3 boundary discussed in Section~\ref{sec:exp:verdicts}.

\paragraph{Stricter gates and content controls.} Three stricter gates were fixed alongside the primary endpoint: gate-C1 requires the high-conflict recall difference to reach $+0.02$; gate-C2 requires the all-object mAP50 to improve by at least $+0.01$; gate-C3 requires SF and AG to be non-inferior within a relative margin of $-1\%$ against the same-budget baseline, with the boundary convention that the threshold is inclusive and the comparison is made on unrounded values, so a row displayed as exactly $-1.00\%$ has to be resolved from its unrounded value. Inference-time content controls complete the design: the learned map can be replaced at test time by an all-ones map or by a fixed seeded permutation of itself (density preserved, placement destroyed). These hooks are part of the module code, not post-hoc edits, and they modify the inference path only, so the controls are comparable to the method arm on the same weights.

\paragraph{Reporting discipline.}
The protocol fixed the narrative disposition for each outcome tier before the arbitration re-test, so that a weak result could not be reframed after the fact. Accordingly, the H5 outcome is reported with fixed wording --- ``statistically significant but below the stricter pre-registered gate thresholds'' --- never as having passed them; negative results (the H1 falsification, the negative distillation result) are reported as first-class findings in Section~\ref{sec:exp}.

\subsection{Computational cost}\label{sec:methods:cost}

CGA's learnable budget is two small convolutional stacks (16-channel hidden width) and one scalar; the Sobel operators are fixed, and the module is attached at a single point at output resolution, so its cost does not scale with backbone depth. Its footprint is selective rather than global: the applied weight is proportional to the conflict map, and at load time the module is an exact no-op. CGA adds $1{,}523$ parameters to the $3{,}562{,}981$-parameter backbone ($+0.043\%$), $2.5\%$ of single-image latency at $507\times507$ ($208.40\pm3.36$ against $203.37\pm1.70$~ms; against the official weights, $208.96\pm0.89$~ms, the difference lies inside the run-to-run spread) and $1.3\%$ at $1034\times769$ ($694.52\pm3.52$ against $685.87\pm1.81$~ms), and no measurable peak memory, since the peak is dominated by backbone activations. Table~\ref{tab:cost} gives the full profile. Both sizes are constrained by the architecture: the patch-embedding stage floors a division by six and the patch-merging stage needs an even feature map, so this implementation accepts only $d \in \{12m+1, 12m+2, 12m+3\}$ with $((d-4)//3+1)$ divisible by four --- which is why the frozen fusion metrics are computed at $507\times507$ (tiled) and generation feeds $1034\times769$ (native, replicate-padded). Neither $512^2$ nor native $1024\times768$ is runnable here, and padding costs $1.1\%$ extra pixels, cropped away before any metric is computed. Backbone FLOPs are not reported, since the profiler counts parameters, latency and peak memory and no FLOP-counting dependency is part of this release; the module's own multiply-accumulates are counted from the convolution output shapes and appear in Table~\ref{tab:cost}.

\paragraph{Measurement settings.} The environment is stated in full because latency figures are only interpretable with it. All timings were taken on a single NVIDIA GeForce RTX 3050 Laptop GPU (4\,GB) under driver 592.27, PyTorch 2.1.0 with its bundled CUDA 11.8 runtime and cuDNN 8.7.0 (\texttt{torch.backends.cudnn.version()} $= 8700$), GPU otherwise idle. Table~\ref{tab:cost} times a block of thirty passes per repeat at \emph{fp32} (no autocast) and reports the mean over three repeats. Because that cannot separate the per-pass distribution from the repeat-to-repeat spread, a second run timed each pass individually with the device synchronised before and after every pass (10 warm-up, $3\times30$ timed, $n = 90$ per cell), confirming the scale: $+2.08\%$ at P50 and $+2.09\%$ at P95 at $507\times507$, and $+1.20\%$ at P50 and $+2.50\%$ at P95 at $1034\times769$. The per-pass protocol carries synchronisation overhead and reports higher absolute levels (CGA $236.6$ against $208.4$~ms at $507\times507$); what is comparable across protocols is the ratio within one protocol.

\begin{table}[t]
\caption{Incremental cost of the arbitration module, measured on one card and one code path with the GPU otherwise idle. ``Params'' counts the full model, so the CGA row shows the module's addition; peak memory is the allocator peak. Latency is the mean of three repeat runs, each ten warm-up passes followed by thirty timed passes with the device synchronised, and the spread across those repeats is given. ``MACs'' counts only the arbitration module's multiply-accumulates, from the convolution output shapes; backbone FLOPs are not counted (see the text). The two sizes are the two this pipeline can run.}
\label{tab:cost}
\scriptsize
\setlength{\tabcolsep}{3.2pt}
\begin{tabular}{@{}llrrrrr@{}}
\toprule
Weights & Size & Params & ms/image & FPS & Peak mem (MB) & MACs (M) \\
\midrule
official & 507$\times$507 & 3{,}562{,}981 & 209.0$\pm$0.9 & 4.79 & 578.25 & --- \\
B0 & 507$\times$507 & 3{,}562{,}981 & 203.4$\pm$1.7 & 4.92 & 578.25 & --- \\
\textbf{CGA} & 507$\times$507 & \textbf{3{,}564{,}504} & \textbf{208.4$\pm$3.4} & \textbf{4.80} & 578.26 & \textbf{378.4} \\
official & 1034$\times$769 & 3{,}562{,}981 & 683.4$\pm$4.6 & 1.46 & 1739.11 & --- \\
B0 & 1034$\times$769 & 3{,}562{,}981 & 685.9$\pm$1.8 & 1.46 & 1739.11 & --- \\
\textbf{CGA} & 1034$\times$769 & \textbf{3{,}564{,}504} & \textbf{694.5$\pm$3.5} & \textbf{1.44} & 1739.12 & \textbf{1170.5} \\
\botrule
\end{tabular}
\end{table}

%%=============================================================
%% §4 Experiments
%% Source: writing/05-experiments.md — unified evaluation protocol, main
%% comparison, cross-seed, hypothesis verdicts, mechanism
%% visualization (2 panel figures), color fidelity, negative
%% results, reproducibility. Tables digit-for-digit.
%%=============================================================
\section{Experiments}\label{sec:exp}

\subsection{Unified evaluation protocol}\label{sec:exp:protocol}

All quantitative evidence is produced by a single evaluation harness, so every row of every table is comparable by construction. Fusion quality is measured on the full 300-image M3FD test split (TarDAL \cite{liu2022tardal}) by the unified evaluator \texttt{eval\_metrics.py} on the luminance channel, reporting six metrics: entropy (EN), spatial frequency (SF), average gradient (AG), mutual information (MI), QABF \cite{xydeas2000objective} and VIF \cite{sheikh2006image}. The same evaluator, channel convention and per-image CSV schema apply to all eight re-implemented baselines and to all three S4Fusion variants.

The eight baselines --- MetaFusion (CVPR'23) \cite{zhao2023metafusion}, CDDFuse (CVPR'23) \cite{zhao2023cddfuse}, EMMA (CVPR'24) \cite{zhao2024equivariant}, DCEvo (CVPR'25) \cite{liu2025dcevo}, W-Mamba (ICME'25) \cite{zhang2025wmamba}, Diff-IF (InfFus'24) \cite{yi2024diffif}, FusionMamba (VI'24) \cite{xie2024fusionmamba}, and the official S4Fusion weights (TIP'25) \cite{ma2025s4fusion} --- are all run under this unified protocol. One protocol note is required for FusionMamba: its VMamba \cite{liu2024vmamba} encoder--decoder requires input side lengths divisible by 32, so 17 of the 300 test images were inferred with replicate padding to a multiple of 32 and cropped back to the original size (\texttt{test\_m3fd\_pad32.py}); the remaining images were inferred at native resolution.

For downstream detection, we reuse the DCEvo detection branch checkpoint (\texttt{DCEvo\_detect\_branch.pt}) with the official M3FD XML$\to$YOLO ground-truth construction, run on the full 300-image test set at input size 640, with a detection counted at IoU $\geq 0.5$. The high-conflict object subset is defined per the protocol from image gradients: for each annotated object, the mean gradient conflict between the infrared and visible source patches is computed as $\min(g_{\mathrm{ir}}, g_{\mathrm{vi}})\cdot(1-\cos\theta)/2$ over the object bounding box, and objects in the top tercile of this conflict score form the high-conflict subset. The same detector, the same ground truth, and the same subset definition are applied to every variant under test.

\paragraph{Data composition and statistical units.} The evaluation population is fixed and reported in full, because every count in this paper is a count over it. The split is the 300-image M3FD manifest (\texttt{dataset/manifests/m3fd\_test.csv}); after the source-pair and minimum-box filters of the object collector it yields $2{,}583$ annotated objects, of which $878$ fall in the high-conflict tercile, and those $878$ are spread over $181$ \emph{images} --- a mean of $4.85$ objects per image, maximum $13$. That nesting is why the detection endpoint is analysed at the image level: objects in one image share its fused output, its detector pass and its local registration. The detector-side environment is inherited from the DCEvo release and is \emph{not} pinned in this paper's artifacts; the fusion-side environment is pinned (\texttt{requirements-brss.txt}, torch 2.1.0+cu118).

Statistical testing follows the protocol: per-image paired Wilcoxon with Holm correction within each family, image-clustered bootstrap and permutation tests as the post-hoc cluster-aware robustness analysis of the detection endpoints, and exact binomial (McNemar-style) tests on recovered/lost objects retained as the pre-specified sensitivity endpoint. The three layers are named in Section~\ref{sec:methods:design} and used with those names everywhere: ``post-hoc cluster-aware robustness analysis'' for the image-clustered reading, ``pre-specified sensitivity endpoint'' for the object-level McNemar-style test, and ``descriptive pooled analysis'' for the across-seed aggregates. All hypotheses, endpoints and thresholds were fixed before the corresponding experiments; where a test fails to reject the text says that no difference was detected rather than that the arms are equivalent.

\subsection{Main comparison}\label{sec:exp:main}

Table~\ref{tab:main} presents the main fusion-quality comparison on M3FD full-300 under the unified evaluation protocol, with the S4Fusion backbone shown as a three-row system: the official released weights, the same-budget retrain (B0), and CGA trained with seed 42.

\begin{table*}[t]
\caption{Fusion quality on M3FD (300 pairs, unified evaluator, Y channel). The ``$\pm$ over'' column states the unit of every dispersion figure in that row, because two different units appear in this table and they are not comparable.}
\label{tab:main}%
\resizebox{\textwidth}{!}{%
\begin{tabular}{@{}llcccccc@{}}
\toprule%
Method & $\pm$ over & EN & SF & AG & MI & QABF & VIF \\
\midrule
S4Fusion (official weights) & 300 images & 6.8926$\pm$0.4159 & 14.9241$\pm$5.4901 & 8.4926$\pm$4.4715 & 4.3364$\pm$1.4031 & 0.5147$\pm$0.0617 & \textbf{0.4628}$\pm$0.0360 \\
MetaFusion (CVPR'23) & 300 images & \textbf{7.2607}$\pm$0.2738 & 13.7397$\pm$4.5319 & 8.2228$\pm$3.8536 & 2.2953$\pm$0.7275 & 0.2189$\pm$0.0415 & 0.2842$\pm$0.0635 \\
CDDFuse (CVPR'23) & 300 images & 6.8996$\pm$0.3864 & 14.7758$\pm$5.8265 & 8.3311$\pm$4.6268 & 3.7648$\pm$0.8611 & 0.4008$\pm$0.0844 & 0.3884$\pm$0.0657 \\
EMMA (CVPR'24) & 300 images & 6.9242$\pm$0.4004 & \textbf{15.2273}$\pm$5.7244 & \textbf{9.0133}$\pm$4.7795 & 3.7391$\pm$0.7523 & 0.3405$\pm$0.0652 & 0.3770$\pm$0.0589 \\
DCEvo (CVPR'25) & 300 images & 6.8301$\pm$0.4369 & 13.7861$\pm$5.1852 & 7.6131$\pm$4.0372 & 4.2834$\pm$0.7955 & 0.4410$\pm$0.0833 & 0.3930$\pm$0.0653 \\
W-Mamba (ICME'25) & 300 images & 6.7083$\pm$0.4029 & 11.8793$\pm$4.9299 & 6.6574$\pm$3.7823 & 4.8307$\pm$0.6497 & 0.3532$\pm$0.0682 & 0.3645$\pm$0.0581 \\
Diff-IF (InfFus'24) & 300 images & 6.7904$\pm$0.4198 & 14.1008$\pm$5.4713 & 7.9352$\pm$4.4519 & \textbf{5.1704}$\pm$0.7274 & 0.4345$\pm$0.0738 & 0.3973$\pm$0.0655 \\
FusionMamba (VI'24) & 300 images & 6.6103$\pm$0.3958 & 11.3349$\pm$4.0729 & 6.2939$\pm$3.2972 & 3.3033$\pm$0.6836 & 0.2296$\pm$0.1198 & 0.3178$\pm$0.0740 \\
S4Fusion (same-budget retrain, B0) & 300 images & 6.8733$\pm$0.3728 & 14.8933$\pm$5.4380 & 8.6801$\pm$4.3975 & 3.2078$\pm$0.8250 & 0.5295$\pm$0.0552 & 0.4535$\pm$0.0509 \\
CGA (ours, s42) & 300 images & 6.8863$\pm$0.3896 & 14.7672$\pm$5.3889 & 8.5173$\pm$4.3712 & 3.9345$\pm$0.8956 & \textbf{0.5449}$\pm$0.0552 & 0.4622$\pm$0.0545 \\
S4Fusion (budget-matched retrain, B0\_e32) & \textbf{3 seeds} & 6.8779$\pm$0.0013 & 14.8577$\pm$0.0026 & 8.6452$\pm$0.0081 & 3.2971$\pm$0.0242 & 0.5342$\pm$0.0002 & 0.4567$\pm$0.0003 \\
CGA (3-seed mean) & \textbf{3 seeds} & 6.8921$\pm$0.0051 & 14.8212$\pm$0.0486 & 8.5465$\pm$0.0265 & 3.9297$\pm$0.0107 & 0.5402$\pm$0.0041 & 0.4611$\pm$0.0010 \\
\botrule%
\end{tabular}}%
\footnotetext{Bold marks the best value per column; the two aggregated rows are marked in the ``$\pm$ over'' column because their dispersion is \emph{across seeds} ($n = 3$), an order of magnitude smaller than the across-image dispersion of an otherwise identical row, and the two must not be compared as if they were the same quantity. B0\_e32 is the \emph{budget-matched} comparator for CGA (both trained for 32 epochs); B0 is the 20-epoch retrain, so CGA$-$B0 bundles the extended schedule while CGA$-$B0\_e32 does not (Section~\ref{sec:exp:verdicts}). The eight baseline rows are re-implemented and evaluated under the unified evaluation protocol of Section~\ref{sec:exp:protocol}; every baseline uses its authors' released checkpoint under its own training recipe and data domain and only the S4Fusion family is retrained in-house, so cross-method comparison carries a provenance caveat. The S4Fusion backbone is additionally shown as the 20-epoch retrain (B0) and with CGA (seed 42, 32-epoch extended schedule; see Sections~\ref{sec:exp} and~\ref{sec:limitations}).}
\end{table*}

Three readings of Table~\ref{tab:main} matter. First, the QABF difference has to be quoted with both the baseline \emph{and} the aggregation level, because matching the levels changes the picture: against the 20-epoch retrain (B0) the seed-42 comparison ($n = 300$, one trained model each side) gives $+0.0155$ ($p = 7.1\times10^{-49}$, $r_{\mathrm{rb}} = +0.98$), while the same convention pooled over three seeds gives $+0.0113$ (95\% CI $[+0.0108, +0.0118]$); against the budget-matched 32-epoch retrain (B0\_e32) the pooled comparison gives $\mathbf{+0.0061}$ (95\% CI $[+0.0056, +0.0066]$). The correct same-level statement about the schedule is therefore $+0.0113$ against $+0.0061$, roughly $1.9\times$, not the $2.5\times$ a single-seed figure against a pooled one would suggest. The comparison against the official weights ($+0.0302$) is context only, since it bundles the retrain's own gain. Second, CGA's VIF is within $-0.0006$ of the official weights and a paired Wilcoxon rejects no difference ($p = 0.5571$); that is not the same as showing equivalence. Third, no single method dominates all six metrics --- EMMA leads SF and AG, Diff-IF leads MI, MetaFusion leads EN --- and no claim is made that the S4Fusion family is the best fusion model overall. Since the CGA arm follows the registered 32-epoch schedule, the CGA$-$B0 deltas should be read as the effect of the recipe rather than of the gate operator alone. The closest conflict-aware methods of Section~\ref{sec:related} are not rows of this table; they are compared on mechanism axes (Table~\ref{tab:positioning}). Figure~\ref{fig:qualitative} complements the table with two scenes against the sources, the official weights, the B0 retrain and three re-implemented baselines.

\begin{figure*}[t]
\centering
\includegraphics[width=\textwidth]{figs/qualitative_01406.png}\\[2pt]
\includegraphics[width=\textwidth]{figs/qualitative_00097.png}
\caption{Qualitative comparison on two M3FD scenes (upper block: conflict-heavy scene 01406, drawn from the six highest-learning conflict samples of the visualization record; lower block: regular scene 00097). Each block has two rows sharing the same eight columns: the upper row is the full frame with the region of interest outlined in red, and the lower row magnifies that region --- $64\times64$ source pixels, the identical window in every column, selected from the infrared source alone as the window of largest gradient energy --- so that target and boundary detail can be compared rather than inferred. All fused columns are produced by the unified evaluation protocol of Section~\ref{sec:exp:protocol}; no image was regenerated for this figure. The complete 300-image sets are released with the code.}
\label{fig:qualitative}
\end{figure*}

FusionMamba is the weakest row of the table across most metrics (EN 6.6103, SF 11.3349, AG 6.2939, QABF 0.2296). We record this honestly with its cause: its released weights are trained on KAIST, a pedestrian-surveillance domain that is out-of-distribution for the road-scene M3FD test set. This row is a domain-transfer observation, not a claim about the FusionMamba architecture; it does not affect the head-of-table conclusions.

\subsection{Cross-seed robustness}\label{sec:exp:seeds}

To rule out single-seed luck, the comparison is replicated over three training seeds (42, 123, 3407), each producing 300 fused image pairs per variant. Table~\ref{tab:perseed} gives the per-seed result, which is the primary reading: each seed is compared against \emph{its own} 32-epoch budget-matched baseline, all three seeds are shown rather than an aggregate, and seed heterogeneity is visible rather than averaged away.

\begin{table}[t]
\caption{Per-seed effect, budget-matched convention (CGA $-$ its own 32-epoch B0), $n = 300$ images per cell, image as the pairing unit, Holm correction over the 18 seed$\times$metric tests. Each cell is the mean per-image difference with its Holm-corrected $p$ in parentheses. This is the primary cross-seed table; Table~\ref{tab:seeds} is its pooled descriptive counterpart and does not carry a replication claim. The last column counts the seeds whose sign agrees with the majority.}
\label{tab:perseed}%
\footnotesize
\setlength{\tabcolsep}{4pt}
\begin{tabular}{@{}lccccl@{}}
\toprule
Metric & s42 & s123 & s3407 & sign \\
\midrule
EN   & $+0.0099$ (1.1e$-$2)  & $+0.0173$ (6.3e$-$6)  & $+0.0154$ (4.3e$-$5)  & 3/3 \\
SF   & $-0.0879$ (1.4e$-$49) & $-0.0255$ (2.0e$-$14) & $+0.0038$ (0.20)      & 2/3 \\
AG   & $-0.1366$ (1.1e$-$49) & $-0.0911$ (3.4e$-$48) & $-0.0686$ (3.9e$-$44) & 3/3 \\
MI   & $+0.6549$ (1.5e$-$48) & $+0.6503$ (1.2e$-$49) & $+0.5928$ (1.3e$-$49) & 3/3 \\
QABF & $+0.0110$ (8.0e$-$47) & $+0.0038$ (1.6e$-$20) & $+0.0035$ (6.3e$-$19) & 3/3 \\
VIF  & $+0.0058$ (1.8e$-$26) & $+0.0037$ (1.2e$-$14) & $+0.0039$ (2.3e$-$16) & 3/3 \\
\botrule
\end{tabular}
\end{table}

Two features of Table~\ref{tab:perseed} matter more than the aggregate. First, the \emph{sign} is stable in all three seeds except SF, which is negative on two and statistically null on the third ($+0.0038$, $p_{\mathrm{Holm}} = 0.20$) --- so the honest statement is not ``SF degrades on 3/3 seeds'' under the matched convention. Second, the \emph{magnitude} is not stable: the matched QABF difference ranges over a factor of three across seeds, and the matched SF cost over a factor of twenty-three once s3407 is included. A direction that replicates while the magnitude moves by a factor of three is a real but weakly identified effect, which is why no single pooled magnitude is attached to a per-seed claim.

Table~\ref{tab:seeds} reports the pooled test for completeness. It is a \emph{descriptive} statistic: it aggregates three already-trained models over one fixed set of 300 images, so its $n = 900$ counts image pairs and not independent training replications, and with $n = 900$ its $p$-values are driven by the pair count rather than by how reproducible the effect is.

\begin{table}[t]
\caption{Pooled per-image paired Wilcoxon (image = unit) with Holm correction, CGA $-$ B0 over 3 seeds $\times$ 300 images ($n = 900$ pairs per metric). \textbf{Descriptive pooled analysis, not a replication test}: the three seeds share the same 300 test images, so this is not 900 independent observations of a training replication. The interval is a 95\% bootstrap over the pooled per-image differences ($B = 2000$, fixed seed); $r_{\mathrm{rb}}$ is the rank-biserial correlation, which is $+1$ when every pair moves the same way and $0$ when the positive and negative ranks balance. Cross-seed reproducibility is carried by Table~\ref{tab:perseed}, not by the $p$-values here.}
\label{tab:seeds}%
\small
\begin{tabular}{@{}lcccr@{}}
\toprule
Metric & Pooled mean $\Delta$ (CGA$-$B0) & 95\% CI & $r_{\mathrm{rb}}$ & $p_{\mathrm{holm}}$ \\
\midrule
EN & $+0.0174$ & $[+0.0114, +0.0235]$ & $+0.318$ & 1.31e$-$16 \\
MI & $+0.7066$ & $[+0.6838, +0.7295]$ & $+0.994$ & 1.68e$-$146 \\
QABF & $+0.0113$ & $[+0.0108, +0.0118]$ & $+0.961$ & 4.71e$-$137 \\
VIF & $+0.0073$ & $[+0.0068, +0.0079]$ & $+0.808$ & 1.26e$-$97 \\
SF & $-0.0713$ & $[-0.0755, -0.0670]$ & $-0.908$ & 1.09e$-$122 \\
AG & $-0.1295$ & $[-0.1335, -0.1254]$ & $-0.997$ & 4.35e$-$147 \\
\botrule
\end{tabular}
\end{table}

Read together, the two tables say this: the information- and fidelity-oriented metrics improve in all three seeds individually under the matched convention, the structural sharpness metrics degrade except that SF is not distinguished from zero on s3407, and no per-seed magnitude may be quoted without its seed. Seed-level tests over $n = 3$ are intentionally omitted, since the minimum attainable Wilcoxon $p$ with three values is 0.25 --- precisely why replication is counted per seed and never inferred from a pooled statistic.

\subsection{Hypothesis verdicts}\label{sec:exp:verdicts}

This section reports the verdicts on the hypotheses, each against the judgment criterion fixed in the protocol before the experiment was run.

\paragraph{H1 --- feature-space $B/C/\Delta$ state modulation improves structural fusion --- criterion met on one seed of three, in opposing directions.}
Under the controlled same-budget ablation (same public S4Fusion starting weights, same budget), the state-modulated variant significantly \emph{degraded} every structural metric relative to B0 on seed 42: SF $-0.0190$, AG $-0.0126$, QABF $-0.0005$, with Holm-corrected $p = 2.90\times10^{-44}$, $8.13\times10^{-40}$, and $5.63\times10^{-27}$ respectively. The protocol's conclusion rule --- ``no structural metric with the state-modulated variant above B0 at Holm $p<0.05$'' --- was met by the negative direction on that seed. Two further seeds trained under the identical budget invalidate the single-seed reading, and they are reported here precisely because the narrative pivoted on this arm: seed 123 meets the same criterion in the \emph{opposite} direction (SF $+0.0328$ and AG $+0.0172$, Holm $p<0.05$, with QABF $-0.0004$ significantly down), and seed 3407 improves AG alone ($+0.0253$) while degrading the QABF median ($-0.0005$). The pre-registered criterion is therefore met on one of three seeds, and the only quantity stable across all three is a uniform small fall in QABF's median. We report the count rather than the verdict: under this parameterization the state-dynamics hypothesis is neither falsified nor supported as a replicated effect, and the pivot to image-space arbitration is defended by the design argument of Section~\ref{sec:methods:problem}, not by this ablation.

The arm is meant to modulate the scan with a \emph{boundary-reliability} estimate, so the protocol carries the control analogous to CGA's: the estimator's output is replaced at inference time by a uniform map and, separately, by a fixed seeded permutation of itself. Neither replacement changes the fused output materially (against the modulated arm, SF moves by $+0.0002$ and $+0.0001$ and AG by $+0.0000$, $p = 0.24$ and $0.60$), and against B0 both reproduce the seed-42 degradation in full (SF $-0.0188$ and $-0.0190$, AG $-0.0126$, QABF $-0.0005$). The effect is therefore not a property of boundary reliability at all: a constant field, carrying no information about the image, and a permuted one, carrying the right values in the wrong places, both produce it. That is the same pattern the CGA controls report (Section~\ref{sec:exp:controls}), and it explains why the H1 verdict changes sign with the seed --- what varies is a global, content-independent change to the scan gains, not a boundary-driven allocation of capacity. We accordingly read this arm as a verdict on the state-capacity account in neither direction.

\paragraph{H5 --- CGA conflict-gated arbitration improves high-conflict detection --- criterion met on one of three seeds, and only that one seed survives the post-hoc cluster-aware analysis.}
The primary endpoint is high-conflict object recall on the full 300-image M3FD detection protocol. Two base conventions are reported below, and they are \emph{not} directly comparable, because they differ in the baseline rather than in the method: the single-baseline convention compares every seed against one shared 20-epoch baseline, whereas the budget-matched convention compares each seed against its own 32-epoch baseline. The second is the fairer comparison, since the method arm trains for 32 epochs; it is not a stricter one.

The post-hoc cluster-aware analysis treats the \emph{image} as the unit of observation, as Section~\ref{sec:methods:design} specifies, because the $878$ high-conflict objects in this subset are carried by only $181$ images (Section~\ref{sec:exp:protocol}) and objects inside one image share its fused output and its detection error. Table~\ref{tab:cluster} gives that analysis under the budget-matched convention. The object-level exact binomial test is reported beside it as a sensitivity analysis, not as the basis of the verdict.

\begin{table}[t]
\caption{Post-hoc cluster-aware robustness analysis of the detection endpoint --- the table from which the detection verdict is read --- under the budget-matched convention, with the image as the unit of observation. ``obj.'' is the object-level discordant count with its exact binomial $p$, i.e.\ the pre-specified sensitivity endpoint and not the basis of the verdict; ``img.'' is the image-level net-sign test; ``flip'' is the exact cluster sign-flip test on the per-image scores; the interval is a percentile bootstrap that resamples the $181$ images carrying the subset with replacement while holding the denominator fixed at the $878$ observed high-conflict objects, so its estimand is the object-weighted high-conflict recall difference of the two fixed arms on the fixed image set ($B = 2000$, fixed seed). The analysis is post-hoc: it was added on 20 September 2026, after the dependence structure was identified and after the re-test outcome it describes had been obtained, so it is a robustness analysis of stored predictions and not a pre-registered test (Table~\ref{tab:timeline}). A seed is counted as supporting the hypothesis only if the image-level test and the bootstrap interval agree. The same $181$ images carry the subset throughout; the remaining images are ties and are excluded from the sign tests.}
\label{tab:cluster}%
\footnotesize
\setlength{\tabcolsep}{4pt}
\begin{tabular}{@{}lcccccc@{}}
\toprule
Seed & base $\to$ CGA & obj. rec/lost ($p$) & img.\ $+/-$ ($p$) & flip $p$ & boot.\ 95\% CI & robust \\
\midrule
s42   & $0.4419\to0.4579$ & 18/4 (4.3e$-$3) & 13/4 (0.049) & 0.021 & $[+0.0046,+0.0296]$ & \textbf{yes} \\
s123  & $0.4476\to0.4544$ & 13/7 (0.26)    & 12/7 (0.36)  & 0.29  & $[-0.0034,+0.0182]$ & no \\
s3407 & $0.4499\to0.4544$ & 11/7 (0.48)    & 11/7 (0.48)  & 0.48  & $[-0.0046,+0.0137]$ & no \\
\botrule
\end{tabular}
\end{table}

The post-hoc cluster-aware result is a marginal image-level result, and it is weaker than the aggregate one: \emph{one} seed of three keeps a positive result under that analysis. Seed 42's image-level test is at the margin ($13$ images up against $4$ down, $p = 0.049$) with a sign-flip $p$ of $0.021$ and a bootstrap interval that excludes zero; on seeds 123 and 3407 the image-level tests do not reject ($p = 0.36$ and $0.48$) and both bootstrap intervals include zero. That is why this paper states the detection finding as a \textbf{seed-specific, non-replicated effect} and not as a gain that replicates across training runs.

Table~\ref{tab:h5e32} repeats the detection endpoint in its object-level form --- the statistic the protocol fixed before the cluster structure was modelled --- and Table~\ref{tab:h5} gives the same test under the single-baseline convention. Both are retained as the pre-specified sensitivity layer of Section~\ref{sec:methods:design}: they are informative in aggregate, and neither is used here as the basis of a replication claim --- that is read from Table~\ref{tab:cluster}.

\begin{table}[t]
\caption{Pre-specified sensitivity endpoint, \emph{single-baseline} convention: high-conflict object recall, CGA vs.\ one shared 20-epoch baseline (B0cmp\_full) for all seeds, at the object level. Not comparable with Table~\ref{tab:h5e32}, and not the basis of the verdict, which is Table~\ref{tab:cluster}. The s123 and s3407 rows are numerically identical because they share this one baseline; against their own 32-epoch baselines the two seeds separate (Table~\ref{tab:h5e32}).}
\label{tab:h5}%
\footnotesize
\setlength{\tabcolsep}{4pt}
\begin{tabular}{@{}lcccr@{}}
\toprule
Seed & Recall $\Delta$ (high-conflict) & recovered / lost & OR [95\% CI] & $p$ (exact binom.) \\
\midrule
s42 & $+0.0194$ (0.4385$\to$0.4579) & 20/3 & 6.67 [1.98, 35.0] & 4.88e$-$04 \\
s123 & $+0.0159$ (0.4385$\to$0.4544) & 18/4 & 4.50 [1.48, 18.3] & 4.34e$-$03 \\
s3407 & $+0.0159$ (0.4385$\to$0.4544) & 18/4 & 4.50 [1.48, 18.3] & 4.34e$-$03 \\
\botrule
\end{tabular}
\end{table}

Table~\ref{tab:h5e32} repeats the same test against the budget-matched baseline. Under this convention the criterion is met on one seed only, and the pooled across-seed test is reported alongside it because the protocol requires both.

\begin{table}[t]
\caption{H5 endpoint in its object-level, pre-specified sensitivity form under the \emph{budget-matched} base convention: CGA vs.\ each seed's own 32-epoch baseline (B0\_e32); read as sensitivity and as a record of the original protocol's own endpoint, never as the basis of the verdict, which is Table~\ref{tab:cluster}; these are the matched-baseline counterparts of Table~\ref{tab:h5}, where the two seeds that share one baseline cannot separate. The odds ratio is recovered/lost, with a 95\% interval from the exact Clopper--Pearson interval on the discordant-pair proportion. "met" is the per-seed criterion ($p<0.05$ and recovered$-$lost $\geq 3$). Not comparable with Table~\ref{tab:h5}.}
\label{tab:h5e32}%
\small
\begin{tabular}{@{}lcccrl@{}}
\toprule
Seed & Base $\to$ CGA ($\Delta$) & recovered / lost & OR [95\% CI] & $p$ & met \\
\midrule
s42 & $0.4419\to0.4579$ ($+0.0159$) & 18/4 & 4.50 [1.48, 18.3] & 4.34e$-$03 & yes \\
s123 & $0.4476\to0.4544$ ($+0.0068$) & 13/7 & 1.86 [0.69, 5.50] & 2.63e$-$01 & no \\
s3407 & $0.4499\to0.4544$ ($+0.0046$) & 11/7 & 1.57 [0.56, 4.78] & 4.81e$-$01 & no \\
\botrule
\end{tabular}
\end{table}

The two conventions agree in direction and disagree in verdict: every seed's object-level delta is positive under both, but the criterion is met on three seeds under the single-baseline convention and on one under the matched one, and the cluster-aware test agrees with that one. We state plainly what that means: \emph{the primary criterion does not replicate}. The pooled across-seed count and the Fisher combination appear only because the protocol requires them, and both are descriptive --- $42$ recovered against $18$ lost ($p = 2.67\times10^{-3}$) and $p = 2.02\times10^{-2}$ respectively. Neither is independent evidence, since the three seeds are evaluated on the same $181$ images with the same detector and ground truth, and neither substitutes for the per-seed count. The identical s123/s3407 rows of the single-baseline table are an artifact of sharing one baseline: against their own 32-epoch baselines those two seeds give $13/7$ and $11/7$, and reading the matched convention as three independent readings or the single-baseline one as two leaves one passing reading (s42 alone) either way.

\paragraph{Stricter gates (C1/C2/C3) --- narrow-margin misses under the single-baseline convention, further misses under the matched one.}
Beyond the primary criterion, the protocol fixed three stricter gates; none is passed in full, under either base convention. Both are reported, since the budget-matched one is the fairer comparison.
\begin{itemize}
\item \textbf{gate-C1} (high-conflict recall $\Delta \geq +0.02$): single-baseline convention, all three seeds FAIL ($+0.0194$ / $+0.0159$ / $+0.0159$, i.e.\ $0.0006$--$0.0041$ below the threshold); budget-matched convention, all three seeds FAIL ($+0.0159$ / $+0.0068$ / $+0.0046$).
\item \textbf{gate-C2} (mAP50 $\Delta \geq +0.01$): single-baseline convention, all three seeds FAIL (s42 $+0.0085 \to 0.7282$; s123 $+0.0028 \to 0.7225$; s3407 $+0.0035 \to 0.7232$, against B0cmp $0.7197$); budget-matched convention, all three seeds FAIL again, with smaller margins (s42 $+0.0084$ vs.\ B0\_e32 $0.7198$; s123 $+0.0012$ vs.\ $0.7213$; s3407 $+0.0016$ vs.\ $0.7216$).
\item \textbf{gate-C3} (SF/AG relative non-inferiority within $-1\%$): single-baseline convention, s123 PASS (SF $-0.39\%$, AG $-0.60\%$) and s3407 PASS (SF $-0.21\%$, AG $-0.36\%$), with s42 the exception on AG; budget-matched convention, SF passes on all three seeds ($-0.5918\%$ / $-0.1716\%$ / $+0.0253\%$) but AG passes only on s3407 ($-1.5780\%$ / $-1.0543\%$ / $-0.7946\%$), so the gate passes on s3407 alone. The s42 AG case under the single-baseline convention is the one boundary this paper has to resolve rather than round, and it resolves against the method: the unrounded relative change is $-1.0043\%$, which is $0.0043$ percentage points below the inclusive $-1\%$ margin, with a one-sided $95\%$ lower bound of $-1.0492\%$ --- so the row that displays as exactly $-1.00\%$ is a fail by a hair's breadth, not a pass. Two boundary conventions are therefore stated explicitly: the margin is inclusive (a value of exactly $-1.0000\%$ would pass) and the comparison is made on unrounded per-image means, because rounding to two decimals would have made this case undecidable. One disclosure follows from auditing these numbers: the two base conventions above were evaluated by two scripts that implement AG differently (a mean of $\sqrt{(g_x^2+g_y^2)/2}$ in the gate script, against the main table's gradient average), so the AG percentages are not on the same scale across the two conventions --- they differ by a roughly constant factor. Every verdict is unchanged by the choice: the single-baseline s42 AG case fails under both implementations ($-1.0043\%$ and $-1.8750\%$ respectively, with one-sided lower bounds $-1.0492\%$ and $-1.9735\%$), and the matched-convention pattern is likewise identical. We report the discrepancy rather than reconcile it, and Table~\ref{tab:main} remains on the \texttt{eval\_metrics.py} scale cited in its caption.
\end{itemize}
The verdict for the grade table is therefore: under the single-baseline convention the primary H5 criterion is met on all three seeds, and under the budget-matched convention on one of three; the three stricter gates are met in full under neither. The paper's wording for the detection gains that do survive is fixed throughout as ``statistically significant but below the stricter pre-registered gate thresholds'', and we never write that CGA ``passed the gates''. Where the two conventions disagree, this paper follows the budget-matched one, because it is the comparison in which method and baseline were given the same training budget, and the recorded consequence of that choice is the conclusion we draw: the gain holds in direction everywhere and in significance on one seed, and it is not a replicated effect.

\paragraph{Two finer statistics for the detector endpoint.} Two finer checks were run afterwards and are reported because they cut in different directions. First, per-image AP is computable offline from the stored predictions, so the mAP50 difference can be given a distribution: under the single-baseline convention a bootstrap over images ($B=2000$, fixed seed) puts the difference at $+0.0161$ with a $95\%$ interval of $[+0.0065,+0.0270]$, which excludes zero but also does not exclude values below the registered $+0.01$ gate, so gate-C2 is not unambiguously met; under the matched convention the same statistic gives $+0.0114$, $+0.0062$ and $+0.0083$, all below the gate. The statistic is not the evaluator's mAP50 --- it averages per-image AP over images, whereas the evaluator averages per-class AP with 101-point interpolation, which is why the levels differ ($0.8048$ and $0.7887$ here against $0.7282$ and $0.7197$ there) --- and both are reported side by side rather than reconciled. Second, the all-object check removes the registration-defined conflict filter: across all $2583$ objects CGA recovers $44$ that the baseline misses and loses $13$ (exact binomial $p = 4.7\times10^{-5}$; recall $0.6384$ against $0.6264$), which directly addresses the selection-bias objection since the high-conflict subset is defined by the same quantity that triggers the mechanism and this difference is not confined to it. Both checks are object-level, so the within-image dependence Table~\ref{tab:cluster} models applies to them too, and both are reported as sensitivity material rather than as the primary reading.

\subsection{Mechanism visualization}\label{sec:exp:vis}

The mechanism-level evidence verifies that CGA acts as designed rather than remaining inert. The global gate scale, which is zero-initialized so that loading the official checkpoint reproduces the baseline exactly, trains to saturation in every seed: $\tanh(\alpha) = 0.99982$ (s42), $0.99871$ (s123) and $0.99889$ (s3407) on the best-validation checkpoints. This reading rules out the self-falsifying scenario in which the gate stays at zero and the measured differences would have to come from some other source; the arbitration path is genuinely open.

\paragraph{The trained selector is effectively binary, and its support is sparse.} The source-selection weight $s$ is continuous by construction, so whether the trained selector behaves like a binary choice is an empirical question about the weights rather than a property of the formula --- which is why the paper calls the mechanism a continuous residual arbitration and reserves \emph{trained near-binary selector} for the measured behaviour. On the seed-42 checkpoint the selector saturates almost everywhere: $P(s \notin [0.1, 0.9]) = 0.9560$, with $P(s \le 0.1) = 0.784$ and $P(s \ge 0.9) = 0.172$. Hardness $\min(s, 1-s)$ falls to $0.0139$ over all pixels and $0.0056$ where the mechanism acts, against $0.1192$ at initialization --- a $20$--$30\times$ sharpening that is what the term \emph{near-binary} rests on. Three further properties are reported with it. The selection is strongly asymmetric, leaning toward the visible source. The support is sparse: $\overline{M} = 0.0447$ and $P(M > 0.5) = 0.33\%$ of pixels. And the learned map departs from the prior it is supervised against: the per-image Spearman correlation between $M$ and $\hat\kappa$ averages $0.4487$, so the map is related to the prior but is not a copy of it. Table~\ref{tab:selector} repeats the measurement on the other two seeds and separates the three quantities the discussion relies on --- how binary the selection is, how dense its support is, and how far the learned map follows the prior; all three move within a narrow band across seeds (saturation $0.935$--$0.956$, hardness $0.0139$--$0.0205$, $\overline{M}$ $0.0301$--$0.0447$, $\rho$ $0.365$--$0.449$), so the near-binary, sparse, prior-related character is a property of the recipe rather than of the seed that carries the headline row.

\begin{table}[t]
\caption{Selector and support statistics on the three trained checkpoints, measured over the full 300-image test split ($2.24\times10^{8}$ pixels per checkpoint). ``hardness'' is $\min(s,1-s)$ averaged over all pixels and over the pixels where the mechanism acts ($M>0.5$); the initialisation reference is $0.1192$. ``sat.'' is $P(s \notin [0.1,0.9])$, and $\overline{M}$, $P(M>0.5)$ and $\rho$ describe the learned support map: its mean, the fraction of pixels it exceeds $0.5$ on, and its mean per-image Spearman correlation with the closed-form prior $\hat\kappa$. The three quantities are kept separate because they move independently: how binary the selection is, how dense the support is, and how far the learned map follows the prior it is supervised against.}
\label{tab:selector}%
\footnotesize
\setlength{\tabcolsep}{4pt}
\begin{tabular}{@{}lcccccc@{}}
\toprule
Seed & $\tanh\alpha$ & hardness (all) & hardness (acts) & sat.\ & $\overline{M}$ / $P(M{>}0.5)$ & $\rho(M,\hat\kappa)$ \\
\midrule
s42   & $0.99982$ & $0.0139$ & $0.0056$ & $0.9560$ & $0.0447$ / $0.32\%$ & $0.4487$ \\
s123  & $0.99871$ & $0.0169$ & $0.0084$ & $0.9464$ & $0.0378$ / $0.26\%$ & $0.4360$ \\
s3407 & $0.99889$ & $0.0205$ & $0.0075$ & $0.9349$ & $0.0301$ / $0.30\%$ & $0.3652$ \\
\botrule
\end{tabular}
\end{table}

At the pixel level the intervention is selective rather than global, and the signature is one of \emph{sparsity} rather than of conflict-correctness. On the six highest learned-conflict samples the mean absolute CGA$-$B0 difference is 0.0125--0.0170 while the per-sample difference-map peaks are 0.5940, 0.5902, 0.5831, 0.5675, 0.6097 and 0.4253, read from the archived arrays. A small mean with a large peak means most of the image is untouched while a small set of pixels is strongly re-assigned; it does not show that those pixels are the conflicting ones, and the control of Section~\ref{sec:exp:controls} tests exactly that and returns a negative. Figures~\ref{fig:panel-a} and~\ref{fig:panel-b} show two panels, each with five columns: the infrared source, the visible luminance, the closed-form conflict prior $\hat\kappa$ (per-image max-normalized, input-derived, no annotation --- the panel labels it ``conflict target''), the learned support map $M$, and the absolute CGA$-$B0 difference map.

\begin{figure}[t]
\centering
\includegraphics[width=0.92\textwidth]{figs/01406_panel.png}
\caption{Mechanism visualization on the highest learned-conflict M3FD sample 01406. Columns: infrared source; visible luminance; closed-form conflict prior $\hat\kappa$ (per-image max-normalized, input-derived, no annotation --- the panel labels it ``conflict target''); learned support map $M$; absolute CGA--B0 difference map. This sample's mean absolute CGA--B0 difference falls inside the 0.0125--0.0170 interval spanned by the six highest-conflict samples (the interval, not the value, is shared across the six), and its difference-map peak is 0.5940: most of the image is untouched while a small set of pixels is strongly re-assigned. Selected conflict panel; full set of six samples in the supplementary material.}
\label{fig:panel-a}
\end{figure}

\begin{figure}[t]
\centering
\includegraphics[width=0.92\textwidth]{figs/01422_panel.png}
\caption{Mechanism visualization on M3FD sample 01422, the second-highest learned-conflict sample, with the same five columns as Fig.~\ref{fig:panel-a}. This sample's mean absolute difference also falls inside the six-sample 0.0125--0.0170 interval, and its per-sample difference-map peak is 0.5902 --- not 0.5940, which belongs to sample 01406 and was mis-attributed to this sample in an earlier version; the value here is read from the archived array for 01422. Selected conflict panel; full set of six samples in the supplementary material.}
\label{fig:panel-b}
\end{figure}

A secondary measurement points the same way from the resolution axis. Under the multi-resolution stress protocol (120 paired images at five resolutions 289--769), the CGA variant shows lower cross-scale drift than B0 at 4 of 5 resolutions, with drift-AUC reduced from 0.8088 to 0.7839 ($\Delta = -0.0249$, $p = 1.41\times10^{-5}$), including the highest resolution tested (769). This measurement comes from the seed-42 arm and is not part of the post-hoc cluster-aware robustness analysis either; it inherits that seed's caveat in full, and it is consistent with the reading that selection rather than averaging carries structure across scales without establishing it.

\subsection{Color fidelity}\label{sec:exp:color}

Because the mechanism moves the fused luminance toward one source at some pixels, its effect on chrominance deserves an explicit check, and the revision extends that check from one seed to all three. Colour metrics are computed on the full 300-image set for every seed against that seed's own budget-matched baseline, with YCbCr reconstruction of the fused output; Table~\ref{tab:color} reports the per-seed differences. The purpose of the three-seed form is that a single checkpoint cannot distinguish a property of the mechanism from a property of one training run, which is the same reason the detection endpoint is reported per seed.

\begin{table}[t]
\caption{Colour-fidelity audit, per seed, against each seed's own 32-epoch baseline (image-paired Wilcoxon, $n = 300$). ``$\Delta$ ($p$)'' per cell; a value in bold is significant at $p<0.05$ and its sign is shared by all three seeds. EOR$\uparrow$ and CIEDE2000$\downarrow$ (lower means closer to the visible reference) are the two directions a mechanism of this kind should improve; the bottom two rows are fidelity descriptors whose degradation would be a cost. The three-seed form is the point: it separates what the mechanism does from what one run happens to do.}
\label{tab:color}%
\footnotesize
\setlength{\tabcolsep}{4pt}
\begin{tabular}{@{}lccc@{}}
\toprule
Metric & s42 & s123 & s3407 \\
\midrule
EOR $\uparrow$ & $+0.0010$ (0.065) & $+0.0007$ (0.96) & $+0.0006$ (0.23) \\
CIEDE2000 $\downarrow$ & $-0.1432$ (0.099) & $\mathbf{-0.1890}$ (1.2e$-$5) & $\mathbf{-0.1290}$ (2.4e$-$3) \\
ECS & $+0.0008$ (1.5e$-$7) & $-0.0006$ (8.3e$-$3) & $-0.0003$ (0.066) \\
OBR & $+0.0002$ (2.9e$-$3) & $-0.0009$ (<1e$-$16) & $-0.0009$ (<1e$-$16) \\
RGB SSIM & $\mathbf{-0.0101}$ (7.9e$-$8) & $\mathbf{-0.0063}$ (0.011) & $\mathbf{-0.0092}$ (1.3e$-$5) \\
Colorfulness & $\mathbf{-0.0034}$ (8.3e$-$11) & $\mathbf{-0.0016}$ (0.023) & $\mathbf{-0.0028}$ (1.9e$-$8) \\
\botrule
\end{tabular}
\end{table}

The three-seed picture is less favourable than the single-seed one and is reported as it is. The two metrics that should improve if the mechanism preserved chrominance behave differently: CIEDE2000 colour difference against the visible reference decreases on all three seeds and significantly on two, a consistent improvement, whereas EOR --- the RGB-space edge-orientation agreement with the sources --- is positive on all three seeds but significant on none ($p = 0.065$, $0.96$, $0.23$, intervals all including zero), so the earlier single-seed reading of this metric as a significant improvement does not survive either the matched baseline or the replication and is withdrawn. The two descriptive quantities, RGB three-channel SSIM and Colorfulness, degrade on all three seeds with significance and track the structural cost seen in the luminance metrics, while ECS and OBR flip sign between seeds and are not interpretable at this sample size. The honest summary is a consistent, mostly significant reduction of colour distance to the visible reference; a consistent, uniformly non-significant edge-orientation change; and a consistent, significant degradation of the two fidelity descriptors, which is the chrominance-side echo of the same structural trade-off rather than a new failure mode.

Two provenance notes belong with this table. The audit is reported at the matched convention because that is the convention the rest of the paper follows; the single-seed figures of the previous version were computed against the 20-epoch shared baseline and are superseded. And the RGB reconstruction path is a separate implementation from the luminance evaluator, so the levels are comparable only within the table --- SF, AG, ECS, RGB SSIM, CIEDE2000 and Colorfulness reproduce the frozen seed-42 values from that path exactly, while EOR and OBR differ from the frozen single-seed table by less than $10^{-2}$ in absolute terms because those two are the metrics the RGB path feeds.

\subsection{Negative results and failure analysis}\label{sec:exp:negative}

This paper's methodological stance is that pre-negative results are findings, and this section reports them without cosmetic repair.

\paragraph{H1 --- the falsification does not replicate (already reported in Section~\ref{sec:exp:verdicts}).} The feature-space $B/C/\Delta$ arm degrades every structural metric on seed 42 (SF $-0.0190$, AG $-0.0126$, QABF $-0.0005$; Holm $p \leq 5.63\times10^{-27}$) but inverts the criterion on seed 123 (SF $+0.0328$, AG $+0.0172$) and leaves seed 3407 with AG alone, with QABF's median down by $0.0005$ on all three. It is reported as a negative result about the \emph{arm}, and as a second instance of the pattern running through this paper: a single-seed verdict in this pipeline does not survive repetition. The arm's own control closes the question of what produced it --- a constant reliability field and a permuted one each reproduce the seed-42 degradation in full --- so the effect is a content-free global change to the scan gains rather than a boundary-driven one, and the ``mechanism-inside-the-SSM'' narrative is retired as unproven rather than refuted.

\paragraph{Distillation --- deterministic failure, reported as a negative result.} The reliability-map distillation arm (teacher = trained CGA seed-42 model with depths 1,2,1; student depths 1,1,1 with \texttt{-{}-use-cga}; output-L1 1.0, feature SmoothL1 0.2, boundary-KL 0.5; 20 epochs, crop 128, seed 42) produced NaN validation loss in epoch 1 and NaN train \emph{and} val losses for epochs 2 through 20; no best checkpoint was ever saved, leaving only an unusable NaN-weights \texttt{\_last.pt}. Per the pre-registration discipline we did not tune hyperparameters and did not rerun, and the component is reported as a first-attempt failure rather than a supported contribution. Two candidate explanations are on record --- the commitment loss diverging in this configuration, possibly the same phenomenon as the late-training NaNs, and an implementation issue in the boundary-KL path on a shallow student --- and neither is tested here.

\paragraph{Late-training NaN instability in the CGA family, and the experiments it blocked.}
Across seeds, the CGA variants exhibit late-epoch validation NaN: seed 123 from epoch 12 and seed 3407 from epoch 13, while in the same replication queue the B0 family completes all 20 epochs without any NaN, and the 32-epoch budget-matched baseline completes all three seeds without NaN as well (96 baseline epochs in total across the two controls). The same instability appears in the seed-42 CGA run, whose best checkpoint is at epoch 26 and whose validation loss goes NaN at epoch 27, so in that seed the selected checkpoint sits one epoch away from divergence. The best checkpoints were saved before NaN onset in every case (s123 best at epoch 11, val 0.3090; s3407 best at epoch 12, val 0.3063), so all reported evaluation numbers come from valid pre-NaN checkpoints, and no result in this paper depends on a post-NaN epoch. We record a mechanistic hypothesis --- the commitment loss (weight 10) plus the separate high learning rate on the module parameters diverging late in training --- as an open item; no hyperparameters were changed and no runs were repeated in this paper, per the same-budget discipline.

That decision has a consequence we state rather than leave implied. Two experiments that would have strengthened the attribution were deliberately \emph{not} run, for the same reason in both cases: the evidence they would be interpreted against did not clear its own gate. A \emph{stable retraining} of the CGA arm would have established whether the effect survives a numerically sound run, and an \emph{orthogonal-factor ablation} --- module only, plus conflict loss, plus commit loss, plus module learning rate, plus extended schedule, all under one budget and one checkpoint rule --- would have isolated which ingredient produces it. Neither was run, because the primary endpoint for this arm is a seed-specific, non-replicated effect (Table~\ref{tab:cluster}); attributing a non-replicating effect to a specific factor would not have changed that finding. The consequence is accepted openly: the method is reported as an \emph{unstable experimental prototype}, the recipe as a bundle, and this paper does not claim to have identified the mechanism of the difference it measures.

\subsection{Content controls of the conflict map}\label{sec:exp:controls}

\paragraph{Both controls, both endpoints, one checkpoint.}
The protocol fixed two inference-time controls for the conflict map: a uniform all-ones map (commitment applied everywhere) and a spatially shuffled map (values preserved, positions destroyed). We ran both at both endpoints, and \emph{with the same checkpoint that produces the main table}, so that the controls and the method are compared like with like --- the earlier, exploratory version of this analysis used the gate-LR variant and is superseded here for exactly that reason. Table~\ref{tab:controls} reports every metric separately, at both statistical layers, because a single summary would hide the fact that these two arms do opposite things to different quantities.

\begin{table}[t]
\caption{Content controls at the detection endpoint. All three arms share one checkpoint and are scored by the identical detector, ground truth and high-conflict object subset. ``obj.'' is the object-level discordant count with its exact binomial $p$, ``img.'' the image-level net-sign test over the $181$ images carrying the subset; ``hc recall'' is the fraction of the $878$ high-conflict objects detected; ``mAP50'' is the evaluator's all-object mAP50 over the $300$ images (the per-image statistic of Section~\ref{sec:exp:verdicts} is given in the text); and ``prec.''/``FP'' are all-object precision and false positives \emph{per image}, all four computed with one estimator for all three arms. The two control rows are deliberately not summarised into a single verdict: they trade different quantities in opposite directions.}
\label{tab:controls}%
\footnotesize
\setlength{\tabcolsep}{4pt}
\begin{tabular}{@{}lccccccc@{}}
\toprule
Map & hc recall & obj.\ rec/lost ($p$) & img.\ $+/-$ ($p$) & mAP50 & prec. & FP/img \\
\midrule
learned ($M$) & $0.4579$ & --- & --- & $\mathbf{0.7282}$ & $0.3816$ & $8.91$ \\
uniform (all-ones) & $\mathbf{0.5103}$ & 77/31 (1.1e$-$5) & 21/24 (0.77) & $0.7161$ & $0.3732$ & $9.26$ \\
shuffled & $0.4567$ & 3/4 (1.00) & 2/3 (1.00) & $\mathbf{0.7305}$ & $0.3812$ & $8.91$ \\
\botrule
\end{tabular}
\end{table}

At the fusion end, per-image paired comparisons against the learned map (300-image M3FD test; five of six metrics, since VIF is not recorded by the evaluator build that produced these runs) show the same split: the all-ones map sits above the learned one on the structural and information metrics (SF by $+2.00$, AG by $+0.35$, MI by $+2.79$, Holm-corrected $p \leq 2.3\times10^{-22}$), with EN slightly lower and QABF not distinguished from it ($-0.0012$, $p = 0.35$); the shuffled map moves every metric by at most $0.44$.

\paragraph{Reading: a trade-off, not a free lunch, and not a placement effect.}
The two controls answer different questions, and only one of them is about this paper's own trigger. The shuffled control is the clean negative: destroying the learned map's arrangement while preserving its values, density and intervention amplitude leaves the detection endpoint where it was ($-0.0011$ recall, $3$ recovered against $4$ lost, $p = 1.00$; image-level $2/3$, $p = 1.00$) and moves the fusion metrics by at most $0.44$. The map's spatial content therefore contributes nothing measurable, and design requirement (R2) of Section~\ref{sec:methods:problem} --- act only where the modalities conflict --- is not supported at either endpoint.

The uniform control is \emph{not} a second negative and not a free variant either; it is a different operating point, and summarising it as ``commitment everywhere does at least as well'' would misreport it in both directions. Committing everywhere raises high-conflict recall substantially and lowers all-object mAP50 below both the learned and the shuffled arm while increasing false positives per image and reducing all-object precision: it buys recall by giving up ranking quality and precision, which is exactly the exchange a deployment decision would have to price. Nor is the recall side as strong as its object-level $p$-value suggests: with the image as the unit the uniform arm is $21$ images up against $24$ down over the $181$ images carrying the subset ($p = 0.77$), and only the exact sign-flip test on the image scores is marginal ($p = 0.0495$); the $1.1\times10^{-5}$ figure is the object-level sensitivity statistic, inflated by counting roughly five objects per image as independent. What the uniform arm establishes is therefore the \emph{operator}: image-space per-pixel selection moves gradient-energy, information and detection quantities whether or not it is confined to conflict pixels. What it does not establish is that the operator is a uniform improvement, and the paper does not claim it is. Two caveats are stated rather than buried: the fusion end carries five of six metrics, and the detection end rests on a single detector; neither affects the direction, which is the same at both endpoints.

\paragraph{Is the commitment free?}
The content controls above test whether the commitment needs the conflict map. A second question is whether it needs to be \emph{learned}: if some parameter-free rule reproduced CGA once it was given the same support, the trained selector would be decoration and what the paper calls a mechanism would be a lookup rule. Table~\ref{tab:free} answers both questions at once, and the answer is negative on every arm.

The rules are defined so that the table can be read without the code. \emph{max} sets the selection weight to $1$ on the modality with the larger local intensity and $0$ on the other; \emph{gradient} does the same by Sobel magnitude, the hard version of the rule the untrained selector already implements. \emph{Global} applies the rule at every pixel, \emph{conflict only} only inside the fixed p85 edge mask used by the earlier exploratory control, and the unqualified rows only inside the mechanism's own learned support $M > 0.5$, which covers $0.33\%$ of pixels. Every arm keeps CGA's architecture, gate and support and changes only the selector, and the two support levels are cross-validated against the same rules, so the influence of the \emph{support} can be separated from that of the \emph{rule}.

Four rules applied through the fixed mask are significantly worse than CGA, the two conflict-confined ones by $0.049$ and $0.156$ recall. Replacing the mask with CGA's own support does not rescue them: hard max and hard gradient inside $M > 0.5$ remain worse than CGA by $0.018$ and $0.022$ ($p = 1.5\times10^{-3}$ and $3.1\times10^{-4}$) and are not distinguished from the same-budget baseline. The sharpest arm is \emph{oracle-init}: CGA's architecture, support and gate unchanged, with only the selector reset to its zero-initialized value. It is worse than CGA by $0.0251$ ($p = 2.7\times10^{-5}$) and slightly below doing nothing at all. What training the selector buys is therefore the difference between a difference and none, in this seed --- and the qualification matters, since every arm here is measured on the seed-42 checkpoint against the shared 20-epoch baseline, i.e.\ against the effect Table~\ref{tab:cluster} finds on that same seed. Table~\ref{tab:free} therefore supports a statement about the \emph{operator} at that operating point and inherits the non-replication recorded above; what it establishes, within that scope, is that the conclusion about the \emph{trigger} does not extend to the \emph{selection}.

\begin{table}[t]
\caption{Parameter-free commitment rules at the detection endpoint. Every arm is built from the same 20-epoch same-budget baseline and scored with the same detector, ground truth and high-conflict object subset as every other detection result in this paper; the rule and support names are defined in the text above. ``Support'' is the set of pixels a rule is allowed to touch: \emph{heuristic} is the fixed p85 edge mask (the earlier exploratory control's mask), \emph{CGA} is the learned map the mechanism itself uses ($M > 0.5$, $0.33\%$ of pixels); the two support levels are cross-validated against the same rules so that the influence of the mask and of the rule can be separated. The last row is the sharpest control: CGA with its selector reset to the zero-initialized value and nothing else changed. $\Delta$ is measured against CGA, so negative means the free rule loses; the same-budget baseline B0 is at recall $0.4385$ throughout, and every arm is measured on the seed-42 checkpoint, so the table inherits that seed's scope. Not comparable with Table~\ref{tab:main}.}
\label{tab:free}%
\footnotesize
\setlength{\tabcolsep}{4pt}
\begin{tabular}{@{}llcccr@{}}
\toprule
Rule & Support & Recall & $\Delta$ vs CGA & rec./lost & $p$ \\
\midrule
CGA (learned selector) & CGA & \textbf{0.4579} & --- & 20/3 & $4.9$e$-$04 \\
\midrule
max, global & heuristic & 0.4237 & $-0.0342$ & 74/104 & 2.9e$-$02 \\
max, conflict only & heuristic & 0.4089 & $-0.0490$ & 54/97 & 5.8e$-$04 \\
gradient, global & heuristic & 0.1697 & $-0.2882$ & 3/256 & 6.3e$-$72 \\
gradient, conflict only & heuristic & 0.3018 & $-0.1560$ & 9/146 & 5.2e$-$33 \\
\midrule
max & CGA & 0.4396 & $-0.0182$ & 4/20 & 1.5e$-$03 \\
gradient & CGA & 0.4362 & $-0.0216$ & 4/23 & 3.1e$-$04 \\
oracle-init (selector at init) & CGA & 0.4328 & $-0.0251$ & 3/25 & 2.7e$-$05 \\
\botrule
\end{tabular}
\end{table}

\subsection{Reproducibility}\label{sec:exp:repro}

The claim made here is deliberately narrower than ``every number can be regenerated''. What the released artifacts support is that every \emph{statistic} can be recomputed from the released per-image and per-object intermediate results without retraining; that the fusion-side pipeline can be re-run end to end from the pinned environment; and that the detection endpoint can be recomputed from the released fused images and predictions, but not guaranteed to reproduce bit-exactly, because the detector library stack is inherited from the DCEvo release and is not pinned in the artifacts accompanying this paper. We state it that way because it is what is true, and the detector environment is recorded as the weakest link rather than left implicit. The version of record for these artifacts is the final review release \texttt{v1.1-final-review-20260921} (21 September 2026), which supersedes the earlier \texttt{v1.0} snapshot and is the version the release manifest hashes; it covers the analyses added after the \texttt{v1.0} freeze --- the budget-matched baselines, the extra feature-space-arm seeds and inference-time controls, the selector-distribution summaries and the cluster-aware re-analysis --- while \texttt{v1.0} remains available and covers the fusion-quality evidence and the content-control detection tables only.

The scripts behind each table: the main comparison table by \texttt{aggregate\_sota.py}; the per-seed and pooled cross-seed tables by \texttt{analyze\_seeds.py} with \texttt{report\_effects.py} for the effect sizes and intervals; the cluster-aware detection analysis of Table~\ref{tab:cluster} by \texttt{cluster\_audit\_final.py}, which prints the object$\to$image$\to$seed mapping, the object- and image-level tests and the bootstrap intervals --- resampling the $181$ images with replacement while holding the denominator fixed at the $878$ high-conflict objects, the convention defined in Section~\ref{sec:methods:design} --- and which asserts the frozen population ($2583$ objects, $878$ high-conflict objects over $181$ images) so that it fails loudly rather than silently analysing a different set; the object-level detection endpoints by \texttt{arb\_full300\_stats.py}, generated from \texttt{gen\_cga\_fused.py} followed by \texttt{arb\_detect\_full300.py} (detector \texttt{ckpt/DCEvo\_detect\_branch.pt}, input size 640, IoU $\geq 0.5$, ground truth from the official M3FD XML annotations); the structural gate by \texttt{struct\_noninferior.py}, with the unrounded boundary values and one-sided intervals of Section~\ref{sec:exp:verdicts} by \texttt{c3\_report.py} and \texttt{c3\_report\_evalmetrics.py}; the latency percentiles by \texttt{profile\_percentiles.py}; the selector distribution by \texttt{export\_select\_stats.py}; the colour audit by \texttt{rebuild\_rgb.py} and \texttt{eval\_metrics\_extended.py}; mechanism visualization by \texttt{cga\_visualize.py}. The seed training chain is orchestrated by \texttt{run\_seeds\_queue\_v2.sh} (marker-file checkpointing) plus \texttt{run\_after\_seeds\_v2.sh}; the FusionMamba baseline via \texttt{test\_m3fd.py}, \texttt{test\_m3fd\_pad32.py} and \texttt{collect\_outputs.py}. The fusion-side environment is pinned in \texttt{requirements-brss.txt} (torch 2.1.0+cu118, Python 3.11) on the single GPU described in Section~\ref{sec:methods:cost}.

\paragraph{What was fixed when: an item-by-item timeline.}
The review of the first submission asked whether the protocol is a full pre-registration. It is not, and the record distinguishes three kinds of claim by when each was fixed relative to the outcome it governs.

\begin{table}[t]
\caption{Provenance timeline. Each row is a record, its date, and the scope of what it fixes. ``Prospective for'' names the outcome the record was fixed \emph{before}. No row is prospective for anything earlier than its own date, and the paper makes no claim that it is.}
\label{tab:timeline}%
\footnotesize
\setlength{\tabcolsep}{4pt}
\begin{tabular}{@{}p{4.45cm}p{1.95cm}p{5.6cm}@{}}
\toprule
Record & Date & Prospective for \\
\midrule
Mechanism protocol (H1 arm, criterion, gates) & 2026-08-13 & the H1 arm and the same-budget ablations that follow \\
Seed-42 CGA checkpoint written & 2026-08-25 & --- (a model artifact, not a criterion) \\
H5 arbitration hypothesis and pass criterion registered & 2026-09-09 & the budget-matched arbitration re-test only \\
Results freeze (main table, baselines, mechanism record) & 2026-09-13 & --- (records results already obtained) \\
Amendment: per-seed Wilcoxon, best-epoch and NaN disclosure & 2026-09-14 & --- (re-analysis of frozen CSVs; no new inference) \\
Pre-registration V1: criterion layering, gate mapping, aggregation rule & 2026-09-14 & the B-3 budget-matched runs, before they were launched \\
Pre-registration V2: independent-readings rule and wording rules & 2026-09-15 & the same runs; wording fixed before outcomes \\
Budget-matched B0\_e32 runs and gate verdicts recorded & 2026-09-15/16 & --- (outcomes) \\
Public release \texttt{v1.0} & 2026-09-15 & --- (a snapshot of the released repository) \\
H1 extra seeds and the content controls & 2026-09-19 & --- (post-hoc additions to the frozen set) \\
Cluster-aware re-analysis (this revision) & 2026-09-20 & --- (post-hoc analysis of stored predictions) \\
Final review release \texttt{v1.1-final-review-20260921} & 2026-09-21 & --- (a snapshot of the released repository; supersedes \texttt{v1.0}) \\
\bottomrule
\end{tabular}
\end{table}

The honest summary of Table~\ref{tab:timeline} is that the \emph{criterion and reporting rules} for the arbitration re-test were fixed before that re-test was run, while the \emph{architecture, training recipe, seed-42 checkpoint and choice of testbed} all pre-date them. The study is a later-stage prospective re-test with a staged lock, not a full pre-registration, and the word is not used here to describe anything larger.

One property of the design deserves emphasis: because the gate is zero-initialized, the trained model can be loaded side-by-side with the official S4Fusion weights, and setting the gate to zero recovers the baseline output exactly ($\tanh(0) = 0$). Any user of the released code can verify that the unmodified official weights are reproduced bit-exactly by the augmented model with the gate closed, and that the reported differences come only from opening it.

%%=============================================================
%% §5 Discussion
%% Source: writing/06-discussion-conclusion.md §6
%%=============================================================
\section{Discussion}\label{sec:discussion}

This section reads the evidence in three registers: what it supports, what it does not, and where the method sits.

\subsection{What the evidence supports}\label{sec:discussion:supports}

The question of this study was narrow by design: at pixels where the modalities disagree, does an explicit image-space per-pixel commitment toward the locally salient modality recover what a pointwise blend loses? Four observations support an answer at the level of mechanism, each attached to the quantity it is about.

First, the gate is genuinely engaged rather than latent: it saturates in every seed (Table~\ref{tab:selector}), which excludes the self-falsification scenario in which the gate stays at zero and any measured difference originates elsewhere. This is the one claim in the paper that holds identically in all three seeds.

Second, the \emph{selector} is what training changes --- the untrained selector is a soft blend, the trained one is near-binary on a small support, and the parameter-free rules that imitate it are all worse than it (Table~\ref{tab:free}) --- so what training buys is located in the selector rather than in the idea of commitment.

Third, the \emph{support} is sparse but its \emph{placement} is not what matters: the learned map acts on a fraction of a percent of pixels, yet a density-preserving permutation of it leaves both endpoints where they were (Section~\ref{sec:exp:controls}). Sparsity is a property of the trained map; selectivity-in-the-right-place is a claim the control refutes, and the two must be kept apart.

Fourth, the \emph{intervention dose} is a real lever: the uniform map touches every pixel rather than $0.3\%$, and it moves high-conflict recall up while moving mAP50, precision and false positives the other way (Table~\ref{tab:controls}). Dose changes the operating point; it does not simply improve it.

The fusion-metric profile is coherent with that reading, within the same scope. Per-seed matched differences (Table~\ref{tab:perseed}) are positive for MI, QABF, VIF and EN in all three seeds and negative for AG in all three, with SF negative on two and undecided on the third. One nuance deserves an explicit reading: relative to the official weights the B0 retrain itself loses MI while gaining QABF, and CGA recovers most of that MI loss on top of B0, so the CGA$-$B0 MI difference is partly a recovery of the retrain's own cost rather than a gain over the official model. Because the evaluator's MI is the sum of its two directed components, the difference decomposes into an infrared-side and a visible-side component ($1.104 \to 1.253$ and $2.104 \to 2.681$), so neither falls and the visible side moves the more of the two --- consistent with the selection asymmetry in Table~\ref{tab:selector} and weighing against the simplest reading, that the difference is only a redistribution of weight between sources, without excluding it. The multi-resolution measurement points the same way (Section~\ref{sec:exp:vis}) but belongs to the seed-42 arm; read together these are mechanism-level observations on one testbed, and the paper does not promote them to an effect-size claim.

\subsection{What the evidence does not support}\label{sec:discussion:notsupport}

Six things the evidence does not support are stated first, because they carry more weight than any of the measurements above. Each is tabulated where it was measured, and the numbers are not repeated here.

\begin{enumerate}
\item \textbf{The detection finding does not replicate.} The budget-matched criterion holds on one seed, and the post-hoc cluster-aware re-analysis leaves that same single seed positive, on a marginal image-level test (Table~\ref{tab:cluster}). The pooled across-seed statistic and the Fisher combination are descriptive only, since the three seeds share the 300 test images. What the paper can say about the detector is a seed-specific, non-replicated effect.
\item \textbf{The property that motivated the mechanism fails its own control.} A density-preserving permutation of the learned map leaves both endpoints where they were (Table~\ref{tab:controls}), so the map's placement does nothing measurable and design requirement (R2) is not supported.
\item \textbf{The uniform arm is a different operating point, not a stronger version of the same result.} It raises high-conflict recall while lowering mAP50 below both the learned and the shuffled arms, raising false positives and lowering precision (Table~\ref{tab:controls}); its object-level $p$-value does not survive clustering. The two arms are not ordered, and the paper does not claim that committing everywhere does at least as well.
\item \textbf{The magnitudes are below the stricter pre-registered gates and the structural cost is real.} C1 and C2 fail on all three seeds under both conventions; C3 passes on two seeds under the single-baseline convention and one under the matched one, with the seed-42 AG case failing by $0.0043$ percentage points on its unrounded value (Section~\ref{sec:exp:verdicts}). Per-seed SF and AG costs are in Table~\ref{tab:perseed}, and the colour audit (Table~\ref{tab:color}) shows the same direction in the two descriptive fidelity descriptors. A mechanism-consistent reading is that selection removes the averaging that partially inflates gradient-energy metrics at boundaries; we report that as the expected price of selection while flagging the reading itself as interpretive.
\item \textbf{Attribution is unresolved.} The arm changes module, losses, learning rate, schedule and checkpoint rule together, and the factor ablation that would separate them was not run (Section~\ref{sec:exp:negative}). The paper measures a recipe; it does not identify the operator as the cause of the difference.
\item \textbf{External validity was not established.} Every number derives from one backbone, one benchmark and one detector, and the one cross-dataset arm in the repository is pre-freeze, uses a different detector and manifest, and reverses direction (Section~\ref{sec:limitations}). The study is scoped to a single-testbed controlled case.
\end{enumerate}

The reliability-map distillation arm is a negative result --- validation NaN from the first epoch, no best checkpoint ever saved --- reported as a failure record, not rehabilitated (Section~\ref{sec:exp:negative}).

\subsection{Positioning}\label{sec:discussion:positioning}

The comparison with the closest conflict-aware methods is defined on mechanism axes, not on gain magnitude, and at the granularity Table~\ref{tab:positioning} records. ARC3Fusion \cite{tian2026adaptive} coordinates consensus--complementarity--conflict in feature space through per-position three-branch Softmax weights, and GatedFusion-Net \cite{brenner2026gatedfusion} learns per-pixel Sigmoid confidence masks; both are continuous, feature-space, soft coordinators. CGA differs on operating space and on the discreteness of the final output: it operates on the two registered source images and its selection sharpens to a near-binary decision after training, but its output is a continuous residual arbitration rather than a discrete pick, so the honest contrast is ``image-space with a trained near-binary selector'' against ``feature-space with a soft weight'', not hard against soft. The trigger axis is no longer part of the claim at all, since Section~\ref{sec:exp:controls} removes it. What survives is the axis distinction, and it is independent of the gain magnitude: a small, non-replicating effect weakens the practical-value claim, not the axis distinction.

The commitment philosophy has a long lineage. The classical select/choose rules \cite{burt1993enhanced,zhang1999categorization,piella2003general} formalized choosing the salient source where sources disagree, and the learning-based multi-focus decision-map line \cite{ma2021endtoend} still makes per-pixel binary choices, but driven by defocus within same-modality pairs. CGA inherits the commitment idea and moves it to cross-modal conflict: the trigger is inter-modal gradient conflict, the selection is learned end-to-end, and it is applied through a zero-initialized gate; the measurement of Section~\ref{sec:exp:vis} then shows that the \emph{learned} part, not the trigger, is what the operator's behaviour rests on.

Finally, the positioning is deliberately that of a single-testbed mechanism audit --- the frame the title carries --- not a comprehensive or class-level superiority claim. EMMA leads SF and AG and Diff-IF leads MI; CGA's leading metric is QABF, with VIF not distinguished from the official weights. S4Fusion is the testbed because its public weights make the comparison controlled, not because it is assumed to be the strongest or most representative backbone, and the conclusions are scoped to S4Fusion, M3FD and the single fixed detector throughout.

%%=============================================================
%% §6 Limitations
%% Source: writing/06-discussion-conclusion.md §7
%%=============================================================
\section{Limitations}\label{sec:limitations}

The following limitations are reported exactly as recorded; none is mitigated by post-hoc argument, and none repeats a result already tabulated.

\begin{enumerate}
\item \textbf{The detection finding does not replicate.} This is the first limitation, not a footnote to one. Under the budget-matched convention the criterion holds on one seed, and under the post-hoc cluster-aware re-analysis the same single seed survives, on an image-level test that is itself marginal. The paper's detector claim is consequently a seed-specific, non-replicated effect, and it should be read as such. The two practical consequences for adopters are that a re-run should be expected to land anywhere in the observed range rather than near the reported point, and that no threshold crossing in this paper should be used as a basis for a deployment decision.

\item \textbf{External validity was not established and was not extended.} The evidence is a single backbone, a single benchmark and a single detector. The extensions that would have supported a broader claim --- an official-manifest re-run of the cross-dataset arm with the frozen checkpoint, a second detector family trained and tuned only on M3FD train/validation, and a controlled $1/2/4$~pixel registration perturbation --- were \emph{not} run. They were not run because the gate that authorizes further performance work (replication of the primary endpoint under the cluster-aware analysis) was not met, and expanding the test bed around a non-replicating endpoint would have produced a broader description of the same non-replication. The study is therefore scoped to an S4Fusion/M3FD/single-detector controlled case, and every class-level phrase has been removed from the abstract, the introduction and the conclusion rather than softened. One residual alternative explanation is also not closed: the intervention shifts local contrast at some pixels, and part of the recall difference could in principle reflect alignment with this detector's preferences rather than improved fusion fidelity. The all-object check, which lies outside the conflict-subset definition, is the evidence against that reading; a second detector would be needed to close it, and no second detector is reported.

\item \textbf{The pre-freeze cross-dataset arm points the other way.} The repository contains an exploratory detection A/B on LLVIP that predates this study's freeze. Its provenance differs from the mainline in three ways: the fused arms were generated on 18 August 2026 from a checkpoint that predates the seed-42 model used for every number here, the detector is a different checkpoint (a DCEvo branch fine-tuned on LLVIP rather than the M3FD branch), and the two runs use different test manifests (1000 and 1050 images). All-object recall is non-positive for the committed arm there ($-0.0080$ and $-0.0004$), i.e.\ the direction is the opposite of the M3FD result. Because the arm is not protocol-comparable it is not evidence against the M3FD numbers; but it is equally not evidence for them out of domain, and on the strict reading of the external-validity gate its direction reversal is what limits the paper's scope. It is reported here rather than omitted.

\item \textbf{Attribution is unresolved.} The CGA arm changes module, losses, learning rate, schedule and checkpoint rule together, so the reported differences are the effect of the recipe. The orthogonal-factor ablation and the stable retraining that would separate these were not run (Section~\ref{sec:exp:negative}), which means the paper measures a recipe and does not identify the operator as its cause. The pre-freeze implementation is labelled \texttt{CGA-v1} so that future work can distinguish it from any factor-resolved successor.

\item \textbf{Training is unstable and the selected checkpoints sit close to divergence.} All three CGA-line seeds show late-epoch validation NaN, while the baseline line does not, and the seed whose checkpoint carries the headline detection row was selected one epoch before its validation loss diverged. The best-validation rule is the only safeguard in this round, and the divergence source is not isolated. The method is therefore reported as an unstable experimental prototype rather than as a finished pipeline; adopters retraining the recipe on their own data should expect to need divergence handling that this study does not provide.

\item \textbf{Structural cost.} Selection removes the averaging that partially inflates gradient-energy metrics at boundaries, so SF and AG fall; the per-seed magnitudes and the gate-C3 boundary case are tabulated in Sections~\ref{sec:exp:seeds} and~\ref{sec:exp:verdicts} and are not repeated here. The reading that this cost is mechanism-inherent rather than training variance is interpretive.

\item \textbf{Registration sensitivity.} The conflict signal $\kappa$ is computed from aligned source pairs; a 1--2 pixel misalignment can turn a shared edge into an apparent gradient-orientation conflict. The experiments use the registered M3FD pairs, no robustness claim is made for misaligned inputs, and the controlled misalignment stress test that would have quantified this was among the extensions not run. Deployments with residual registration error should verify alignment upstream of the gate.

\item \textbf{Per-pixel independent gating.} The selection is made independently at each pixel from local gradient conflict, with no spatial regularization on the selection field. The observed concentration of modifications is an empirical property of the trained map, not an architectural guarantee; spatially regularized selection fields remain untested.

\item \textbf{Budget and convergence.} The budget-matched control had not converged at its budget in any seed (its best validation epoch is the final epoch), so the matched comparison answers ``at equal budget'' rather than ``after convergence'', and a longer baseline budget could only shrink the differences further. The single-seed deltas of a few thousandths also carry run-to-run noise of the same order, since two runs of one arm with one seed do not reproduce each other to better than a few thousandths of validation loss; such deltas must not be read as mechanism effects, which is why the paper's claims rest on per-seed counts and cluster-aware tests rather than on any single small number.

\item \textbf{Reproduction of the detection pass is not guaranteed.} The fusion-side environment is pinned, but the detector-side library stack is inherited from the DCEvo release and is not pinned in the artifacts accompanying this paper. The statistics can be recomputed from the released intermediate results; the detector pass itself is the weakest link, and the paper's reproducibility claim is scoped accordingly in Section~\ref{sec:exp:repro}.
\end{enumerate}

%%=============================================================
%% §7 Conclusion
%% Source: writing/06-discussion-conclusion.md §8
%%=============================================================
\section{Conclusion}\label{sec:conclusion}

This study set out to locate the residual failure of a scan-and-blend selective state-space fusion backbone at modal boundaries, on one testbed. The first arm did not settle the question: feature-space modulation of the state dynamics degrades the structural metrics on one seed under a same-budget protocol, meets the same criterion in the opposite direction on a second, and lowers QABF's median on all three, and the arm's own control shows why the sign moves --- a constant reliability map and a permuted one each reproduce the seed-42 degradation in full. The arm tested next intervenes at conflicting pixels instead, on the design argument rather than on that verdict.

Four propositions are what this paper is prepared to state, each bounded by the evidence named with it.

\begin{enumerate}
\item \textbf{A zero-initialized image-space selection adapter can be added to a released fusion checkpoint without perturbing it.} With the gate closed the augmented model reproduces the official output exactly, and after training the gate saturates in every seed, so the adapter is active rather than latent. This claim is about the implementation, and it is the one that holds identically in all three seeds.
\item \textbf{The conflict-triggered restriction that motivated the design is not supported.} Permuting the learned support map changes neither the fusion metrics nor the detection endpoint (Section~\ref{sec:exp:controls}), so the mechanism is supported, if at all, as a learned operator rather than as a conflict-triggered one.
\item \textbf{The learned selector is not reproduced by any parameter-free rule tested.} Seven rules are all worse than the trained selector on the seed where the effect appears (Table~\ref{tab:free}); what training the selector buys is the difference between a difference and none.
\item \textbf{The detection benefit does not replicate.} The budget-matched criterion holds on one seed, a post-hoc cluster-aware re-analysis leaves the same single seed positive on a marginal image-level test, every other magnitude is below its stricter pre-registered gate, and the whole study rests on one detector and one dataset. The detector claim is a seed-specific, non-replicated effect within a single-testbed controlled case.
\end{enumerate}

Two extensions that would have widened the claim --- a factor-resolved ablation of the training recipe, and an external-validity programme on a second dataset, a second detector and controlled misregistration --- were not run, because the gate that authorizes further performance work did not open; their absence is why this paper measures a recipe rather than an operator, and why it is scoped to one test bed. Two properties of the pipeline bound what any individual number means: the baseline had not converged at the matched budget, and two runs of one arm with one seed do not reproduce each other to better than a few thousandths of validation loss. Within these boundaries what remains is a mechanism that is implemented and active but not a reliable source of gain, a set of negative results that are legible because the protocol made them so, and a protocol --- criterion layering, both base conventions reported as non-comparable, controls at both endpoints, cluster-aware inference, replication counted per seed --- that is the part of this work most directly reusable.

Whether spatially regularized selection fields, a stabilized training recipe, factor-resolved attribution or cross-dataset validation make the mechanism reliable remains open; they are questions, not promises, and the negative results above are why.

%%=============================================================
%% Declarations
%%=============================================================
\backmatter

\bmhead{Data availability}The M3FD benchmark used in this study is publicly available from the TarDAL project \cite{liu2022tardal}; the source datasets are not redistributed here and must be obtained from their original providers. No new data were created.

The public release accompanying this article is the final review release, tagged \texttt{v1.1-final-review-20260921} at \url{https://github.com/JassonG-xt/s4fusion-cga/releases/tag/v1.1-final-review-20260921} and frozen on \textbf{21 September 2026}. It is the version intended to cover every number reported here: the per-image fusion-quality CSVs behind the main comparison, the per-seed and pooled cross-seed tables and the colour audit; the per-image and per-object detection tables behind the learned/uniform/shuffled content-control analysis; the budget-matched 32-epoch baselines and their per-seed comparisons; the additional feature-space-arm seeds and their inference-time controls; the selector-distribution summaries; the object$\to$image$\to$seed mapping; the manifest and checkpoint hashes; and the cluster-aware audit script whose output reproduces Tables~\ref{tab:cluster} and~\ref{tab:controls} from the stored per-object detection vectors. It supersedes the earlier tag \texttt{v1.0} (15 September 2026), which remains available and should be read as covering the fusion-quality evidence and the content-control detection tables only: \texttt{v1.0} was frozen before the budget-matched runs (16 September 2026), the additional feature-space-arm seeds and the inference-time controls (19 September 2026) and the cluster-aware re-analysis (20 September 2026) were produced, and it is therefore not a snapshot of every number in this article.

\bmhead{Code availability}The scripts named in Section~\ref{sec:exp:repro} are released in the same final review release (\texttt{v1.1-final-review-20260921}) described under Data availability, under the licences recorded in its \texttt{LICENSE} file (MIT for code, CC-BY-4.0 for documentation, ledgers and tabular results). One scope limit is stated rather than implied: the detector-side library stack is inherited from the DCEvo release and is \emph{not} pinned in the released artifacts, so the detection endpoint is reproducible by recomputing the statistics from the released intermediate predictions, not by re-running the detector and expecting bit-identical output. The fusion-side environment \emph{is} pinned (\texttt{requirements-brss.txt}, torch 2.1.0+cu118, Python 3.11), and the detector-side environment as observed is recorded in the release's environment notes together with the commands that regenerate each table and figure.

The protocol records are archived rather than summarised. The frozen results list is \texttt{GATE1\_FREEZE.md} (13 September 2026) and the pre-registered decision rules are \texttt{GATE1\_PREREG\_B3\_20260914.md} (14 September 2026), both in the release's \texttt{ledgers/} area; the wording-and-evidence amendment of 15 September 2026 (\texttt{GATE1\_PREREG\_B3\_V2\_20260915.md}), the matched-budget record of 15--16 September 2026 (\texttt{GATE1\_UNFREEZE\_20260915.md}) and the amendment of 14 September 2026 (\texttt{GATE1\_AMEND\_20260914.md}) are retained in the author's working repository and are available to the editor and the reviewers on request. What those records do and do not fix is set out item by item in Table~\ref{tab:timeline}: the criterion, the gate thresholds and the reporting rules for the arbitration re-test were fixed before that re-test was run, while the architecture, the training recipe, the seed-42 checkpoint and the choice of testbed pre-date them, and the cluster-aware re-analysis post-dates the outcome it describes. This study is therefore a later-stage prospective re-test with a staged lock, not a full pre-registration, and the ordering claimed here --- each criterion fixed before the outcome it governs --- is checkable against the archived files' dates and their internal record of the runs.

\bmhead{Author contribution}Xiaotong Gao: conceptualization, method design, experiments, analysis, writing, final verification of every reported number and statement, and responsibility for the content and integrity of the article.

\bmhead{Funding}No funding was received for this work.

\bmhead{Competing interests}The author declares no competing interests.

\bmhead{Ethics approval}Not applicable; this study uses a publicly released benchmark and contains no human participants or personal data.

\bmhead{Use of AI tools}In preparing this manuscript, the author used AI-based assistance for six tasks: literature verification (cross-checking reference identities, the DOI metadata of the two closest published neighbours, and the closest-neighbour characterizations cited in Section~\ref{sec:related}, where the concurrent preprint was verified at abstract level only and Table~\ref{tab:positioning} records that difference); drafting assistance for the research design and protocol text, including the layering of the primary criterion, the stricter gates C1--C3, the decision rules and the wording rules for reporting each outcome tier; the decision to reframe the study as a controlled falsification and mechanism audit after the arbitration arm failed to replicate; statistical re-analysis scripting, including the cluster-aware object$\to$image$\to$seed analysis, the non-inferiority boundary report and the effect-size and interval computations; implementation assistance on analysis and figure code; and language polishing.

The design-related participation is disclosed explicitly because it is the part most easily under-reported. Three boundaries apply to it. First, every criterion, threshold and decision rule was reviewed and confirmed by the author and frozen before the corresponding experiment was run, and the frozen records are those named under Code availability; the post-hoc cluster-aware re-analysis of Section~\ref{sec:methods:design} is labelled as post-hoc precisely because it was not part of that frozen set. Second, the \emph{constraints} that produced the negative findings reported here were fixed by those records rather than chosen at write-up time: the per-seed replication requirement, the inclusive $-1\%$ non-inferiority margin, the requirement to report both base conventions as non-comparable, and the requirement to report failed gates were all registered beforehand. Third, no AI system retrained a model, ran the detector, or generated a number that appears in this article; all experimental numbers were produced by the frozen training and evaluation pipeline, and every re-analysis reported in the revision was computed by a script archived with the submission and re-run against the stored intermediate results. Every AI-assisted statement and number was verified against those records by the author, who reviewed the final text in full, approved each numeric claim and each scoping statement, and takes full responsibility for the content and integrity of the article.

%%=============================================================
%% Bibliography (sn-mathphys-num, numeric citations)
%%=============================================================
\bibliography{refs}

\end{document}
